%
% Tex input file for CGNS SIDS-to-HDF File Mapping Manual
%
% To generate a DVI file, then a PostScript file named filemap_hdf.ps:
%
%     latex filemap_hdf.tex
%     dvips filemap_hdf.dvi -o
%
% To generate a PDF file named filemap_hdf.pdf:
%
%     pdflatex filemap_hdf.tex


\documentclass[twoside,fleqn]{article}

%
% Packages
%
\usepackage{indentfirst}		% indent first paragraph in sections
\usepackage[tbtags]{amsmath}		% use ams math package
\usepackage{graphicx}			% for including eps files
\usepackage{lscape}			% for landscape stuff
\usepackage[bf]{caption}		% for more flexible captions
   \setlength{\captionmargin}{18pt}	% caption margins
\usepackage{array}			% use new array and table package
   \setlength{\extrarowheight}{2pt}	% increase row heights in tables
\usepackage{longtable}			% use longtable package
\usepackage{alltt}			% verbatim input with font changes
\usepackage{natbib}			% for author-year citations
   \setlength{\bibhang}{0pt}		% don't use hanging indent format
\usepackage[normalem]{ulem}		% allow underlining
\usepackage{calc}			% calculation package
\usepackage{ifthen}			% control structures
\usepackage{fancyhdr}			% fancy headers and footers
   \fancyhf{}
   \renewcommand{\headrulewidth}{0pt}
   \renewcommand{\footrulewidth}{0pt}
\usepackage{titling}			% more title page control
\usepackage{xspace}			% common-sense spacing after text macro
\usepackage{mdwlist}			% more flexible description lists
\usepackage{epic}			% extended picture environment
%\usepackage[pdftex]{color}		% colors
\usepackage{color}			% colors
   \definecolor{subcolor}{rgb}{0.809,0.852,1.0}
\usepackage{pdfcolmk}			% for colors across page breaks
%\usepackage[T1]{fontenc}		% for better-looking "_" and tt "{"
\usepackage{ae}				% used instead of above for PDF output
\usepackage{hyperref}			% for hypertext links in PDF
   \hypersetup{letterpaper,plainpages=false,
               pdftitle={CGNS SIDS-to-HDF File Mapping Manual},
               pdfauthor={CGNS Project Group},
               colorlinks,
               linkcolor=blue,citecolor=blue,filecolor=red,pagecolor=blue,
               urlcolor=red}
   \renewcommand{\sectionautorefname}{Section}
   \renewcommand{\subsectionautorefname}{Section}
   \renewcommand{\subsubsectionautorefname}{Section}
   \renewcommand{\paragraphautorefname}{Section}

%
% Page layout
%
\normalsize
\setlength{\oddsidemargin}{0.5in}	% left margins-1.0in for odd/even pages
\setlength{\evensidemargin}{0.0in}
\setlength{\textwidth}{6.0in}		% text width
\setlength{\topmargin}{-0.5in}		% top margin-1.0in
\setlength{\headheight}{0.25in}		% header space parameters
\setlength{\headsep}{0.5in}
\setlength{\textheight}{8.5in}		% text height
\setlength{\topskip}{\baselineskip}	% dist from top of body to 1st baseline
\setlength{\footskip}{0.75in}		% dist from bottom of body to footer
\raggedbottom				% to avoid stretching vertical space

\pagestyle{fancy}			% using fancyhdr package
\fancypagestyle{plain}{%		% page # in corner for plain style
   \fancyhf{}
   \fancyfoot[LE,RO]{\bfseries \thepage}}
\addtocounter{secnumdepth}{1}		% number \paragraph's
\addtocounter{tocdepth}{1}		%    and include them in TOC

%
% Misc spacing parameters
%
\setlength{\parskip}{0.5\baselineskip}	% blank space between paragraphs
\setlength{\doublerulesep}{0pt}		% for wider \hlines
\setlength{\fboxsep}{5pt}		% for more margin in boxes
\newlength{\saveparindent}		% save for use where LaTeX changes them
\newlength{\saveparskip}		% save for use where LaTeX changes them
\newlength{\savebaselineskip}
\setlength{\saveparindent}{\parindent}
\setlength{\saveparskip}{\parskip}
\setlength{\savebaselineskip}{\baselineskip}
\newlength{\tmplength}			% for temp use wherever needed
\newlength{\Pwidth}			% for use in tables
\newlength{\Pwidtha}

%
% User-defined environments
%
% "Definition" list with variable-length entries. Requires
% mdwlist package. (See LaTeX Companion, p 64, and mdwlist documentation.)
\newenvironment{Ventry}[1]%
   {\settowidth{\tmplength}{#1\hspace{1em}}%
    \begin{basedescript}{%
           \desclabelwidth{\tmplength}%
           \desclabelstyle{\nextlinelabel}%
           \renewcommand{\makelabel}[1]{##1\hfil}}}
   {\end{basedescript}}
% "Definition" list with variable-length entries, indented. Requires
% mdwlist package. (See LaTeX Companion, p 64, and mdwlist documentation.)
\newenvironment{Ventryi}[1]%
   {\settowidth{\tmplength}{\hspace{\saveparindent}#1\hspace{1em}}%
    \begin{basedescript}{%
           \desclabelwidth{\tmplength}%
           \desclabelstyle{\nextlinelabel}%
           \renewcommand{\makelabel}[1]{\hspace{\saveparindent}##1\hfil}}}
   {\end{basedescript}}
% Indented definition list with variable-length entries, compact.
% (see LaTeX Companion, p 64)
%\newenvironment{Ventryic}[1]%
%   {\begin{list}{}{\renewcommand{\makelabel}[1]{\hspace{\saveparindent}##1\hfil}%
%      \settowidth{\labelwidth}{\hspace{\saveparindent}#1\hspace{1em}}%
%      \setlength{\leftmargin}{\labelwidth+\labelsep}
%      \setlength{\parsep}{0in}%
%      \setlength{\itemsep}{0in}}}%
%   {\end{list}}
% Indented definition list with variable-length entries, compact. Requires
% mdwlist package. (See LaTeX Companion, p 64, and mdwlist documentation.)
\newenvironment{Ventryic}[1]%
   {\settowidth{\tmplength}{\hspace{\saveparindent}#1\hspace{1em}}%
    \begin{basedescript}{%
           \desclabelwidth{\tmplength}%
           \desclabelstyle{\nextlinelabel}%
           \renewcommand{\makelabel}[1]{\hspace{\saveparindent}##1\hfil}
           \setlength{\topsep}{0in}%
           \setlength{\parsep}{0in}%
           \setlength{\itemsep}{0in}}}
   {\end{basedescript}}
% Definition list with variable-length entries, compact. Requires
% mdwlist package. (See LaTeX Companion, p 64, and mdwlist documentation.)
\newenvironment{Ventryc}[1]%
   {\settowidth{\tmplength}{#1\hspace{1em}}%
    \begin{basedescript}{%
           \desclabelwidth{\tmplength}%
           \desclabelstyle{\nextlinelabel}%
           \renewcommand{\makelabel}[1]{##1\hfil}
           \setlength{\topsep}{0in}%
           \setlength{\parsep}{0in}%
           \setlength{\itemsep}{0in}}}
   {\end{basedescript}}

% "changes" environment, to identify changed code by color.  The color is
% given by the argument, with a default of red.  When the environment is
% exited, the color is changed to black.
\newenvironment{changes}[1][red]%
   {\color{#1}}
   {\color{black}}

% Adapted from comp.text.tex post by Keith Reckdahl
% (reckdahl@am-sparc7.Stanford.EDU)
% Indents text from the left margin by current value of the list
% length \leftmargin
\newenvironment{indleft}%
   {\begin{list}{}
      {\setlength{\topsep}{0pt}%
       \setlength{\listparindent}{0pt}%
       \setlength{\itemindent}{0pt}%
       \setlength{\parsep}{\parskip}%
      }%
      \item[]}%
   {\end{list}}
% As above, but surrounding an alltt environment
\newenvironment{indlefttt}%
   {\begin{list}{}
      {\setlength{\topsep}{0pt}%
       \setlength{\listparindent}{0pt}%
       \setlength{\itemindent}{0pt}%
       \setlength{\parsep}{\parskip}%
      }%
      \item[]%
      \begin{alltt}}%
   {\end{alltt}\end{list}}
% As above, but also with a tabbing environment
\newenvironment{tabindlefttt}%
   {\begin{list}{}
      {\setlength{\topsep}{0pt}%
       \setlength{\listparindent}{0pt}%
       \setlength{\itemindent}{0pt}%
       \setlength{\parsep}{\parskip}%
      }%
      \item[]%
      \begin{alltt}\begin{tabbing}}%
   {\end{tabbing}\end{alltt}\end{list}}

%
% User-defined commands
%

% Redefine \paragraph as a display heading
\makeatletter
\renewcommand{\paragraph}{\@startsection
   {paragraph}%			% heading name
   {4}%				% heading level
   {0in}%			% indentation of heading
   {0.5\baselineskip}%		% space before heading
   {0.1\baselineskip}%		% space after heading
   {\normalfont\normalsize}}%	% style
\makeatother

% Degree symbol
\newcommand{\degree}[1][]{\ensuremath{{#1}^\circ}\xspace}

% Degree Rankine symbol
\newcommand{\degR}{\textsuperscript{$\circ$}\textnormal{R}\xspace}

\newcommand{\HDF}{HDF5\xspace}
% --- ADFH should have a ``name'', like MLL
% --- SLL (Storage level library), LLL (Low level library),
% --- PLL (Physical level library), NLL (Node), BLL (Below), FLL (File)
% --- GLL (Generic), KLL (Kernel), CLL (Core), WLL (Wouaw), ZLL (Zoo)
\newcommand{\ADFH}{SLL\xspace}
\newcommand{\SLL}{SLL\xspace}

% ############################################################################
\newcommand{\HDFend}{%
\marginpar[\begin{flushright}%
\vskip -1.5em%
$\bigtriangleup$$\bigtriangleup$$\bigtriangleup$%
\end{flushright}
]{%
\begin{flushright}%
\vskip -1.5em%
$\bigtriangleup$$\bigtriangleup$$\bigtriangleup$%
\end{flushright}
}}
%
\newcommand{\HDFstart}{%
\marginpar[\begin{flushright}%
\vskip 0.5em%
$\bigtriangledown$$\bigtriangledown$$\bigtriangledown$%
\end{flushright}
]{%
\begin{flushleft}%
\vskip 0.5em%
$\bigtriangledown$$\bigtriangledown$$\bigtriangledown$%
\end{flushleft}%
}}
%
\newcommand{\HDFline}{%
\marginpar[\begin{flushright}%
\vskip 1em%
{\LARGE $\triangleright$$\triangleright$$\triangleright$}%
\end{flushright}
]{%
\begin{flushleft}%
\vskip 1em%
{\LARGE $\triangleleft$$\triangleleft$$\triangleleft$}%
\end{flushleft}%
}}
% ############################################################################


% Re-def longtable's caption command to use \captionlabelfont from caption pkg
% (basis lifted from longtable.sty)
\makeatletter
\renewcommand{\LT@makecaption}[3]{%
  \LT@mcol\LT@cols c{\hbox to\z@{\hss\parbox[t]\LTcapwidth{%
    \sbox\@tempboxa{{\captionlabelfont #1{#2: }}#3}%
    \ifdim\wd\@tempboxa>\hsize
      {\captionlabelfont #1{#2: }}#3%
    \else
      \hbox to\hsize{\hfil\box\@tempboxa\hfil}%
    \fi
    \endgraf\vskip\baselineskip}%
  \hss}}}
\makeatother

% Make next page odd, with preceding blank page empty (LaTeX Companion, p 93)
\newcommand{\clearemptydoublepage}{\newpage{\pagestyle{empty}\cleardoublepage}}

% Text subscripts, analogous to \textsuperscript command (from comp.text.tex
% post by rf@cl.cam.ac.uk (Robin Fairbairns))
\makeatletter
\DeclareRobustCommand*\tsub[1]{%
  \@tsub{\selectfont#1}}
\def\@tsub#1{%
  {\m@th\ensuremath{_{\mbox{\fontsize\sf@size\z@#1}}}}}
\makeatother

% Text superscripts, but shorter
\newcommand{\tsup}[1]{\textsuperscript{#1}}

% Shortcuts for specific fonts
\newcommand{\bold}[1]{{\normalfont\textbf{#1}}}  % Bold
\newcommand{\ital}[1]{{\normalfont\textit{#1}}}  % Italic
\newcommand{\key} [1]{{\normalfont\texttt{#1}}}  % Fixed font for keywords
\newcommand{\fort}[1]{{\normalfont\texttt{#1}}}  % Fixed font for Fortran stuff

% Shortcuts for math character sizes
\newcommand{\Dstyle}{\displaystyle}
\newcommand{\Tstyle}{\textstyle}
\newcommand{\Sstyle}{\scriptstyle}

% Use section titles as marks (i.e., in headers) instead of subsection titles
\renewcommand{\sectionmark}[1]{\markright{\thesection\ \ #1}}
\renewcommand{\subsectionmark}[1]{}

% "Better" treatment of headings for Appendices  (LaTeX Companion, pp 29,30)
%\renewcommand{\appendix}{%
%   \newcommand{\app}{%
%      \secdef\Appendix\sAppendix}%
%   \newcommand{\subapp}{%
%      \secdef\subAppendix\subsAppendix}%
%   \newcounter{app}%
%   \newcounter{subapp}[app]%
%   \renewcommand{\theapp}{\Alph{app}}%
%   \renewcommand{\thesubapp}{\theapp.\arabic{subapp}}%
%   \renewcommand{\sectionmark}[1]{\markright{\theapp\ \ ##1}}
% }
%\newcommand{\Appendix}[2][?]{%   Complex form
%   \refstepcounter{app}%
%   \addcontentsline{toc}{section}%
%      {\protect\bfseries{\appendixname~\theapp.\ \ } #1}%
%   {\noindent\Large\bfseries\appendixname\ \theapp.\ \ #2\par}%
%   \sectionmark{#1}\vspace{\baselineskip}}
%\newcommand{\sAppendix}[1]{%   Starred form
%   \refstepcounter{app}%
%   {\noindent\Large\bfseries\appendixname\ \theapp.\ \ #1\par}%
%   \vspace{\baselineskip}}
%
%\newcommand{\subAppendix}[2][?]{%   Complex form
%   \refstepcounter{subapp}%
%   \addcontentsline{toc}{subsection}%
%       {\thesubapp\ \ #1}%
%   \vspace{\baselineskip}
%   {\noindent\large\bfseries\thesubapp\ \ #2\par}%
%   \vspace{\baselineskip}}
%\newcommand{\subsAppendix}[1]{%   Starred form
%   \refstepcounter{subapp}%
%   \vspace{\baselineskip}
%   {\noindent\large\bfseries\thesubapp\ \ #1\par}%
%   \vspace{\baselineskip}}

\begin{document}

\pagenumbering{roman}

%\fancyhead[LE]{\bfseries SIDS-to-HDF File Mapping Manual}
%\fancyhead[RO]{\bfseries \rightmark}
\fancyfoot[LE,RO]{\bfseries \thepage}

%\setlength{\droptitle}{1.0in}
\pretitle{\begin{flushleft}\LARGE%
          \includegraphics[width=2.0in]{../logo/cgns_bw}\\*[0.25in]}
\posttitle{\par\end{flushleft}\vskip 1.0em}
\title{{\bfseries CFD General Notation System\\
SIDS-to-HDF File Mapping Manual}\\*[0.25in]
{\Large Document Version 1.1\\
CGNS Version 2.5}}
\author{}
\date{}
\maketitle
\thispagestyle{empty}

%\title{{\LARGE \textbf{CGNS}}\\*[0.5in]
%The CFD General Notation System\\
%SIDS-to-HDF File Mapping Manual\footnote
%{This is Version \textcolor{red}{1.3 of this document, released ?? ??? 2002}.
%The current PDF version was derived from the MS Word version of the
%document, created by Boeing.}}
%\author{CGNS Project Group}
%\date{}
%\maketitle
%\thispagestyle{empty}

\clearemptydoublepage
\setlength{\parskip}{0ex}		% remove blank space between paragraphs
\thispagestyle{plain}
\tableofcontents
\listoffigures
%\listoftables
\setlength{\parskip}{\saveparskip}	% put it back

\clearemptydoublepage
\pagenumbering{arabic}

\fancyhead[LE]{\bfseries SIDS-to-HDF File Mapping Manual}
\fancyhead[RO]{\bfseries \rightmark}

\section{Brief Description of CGNS}
\label{s:brief}
\thispagestyle{plain}

The CFD General Notation System (CGNS) is a set of standards, together
with software implementing those standards, for the recording of data
associated with Computational Fluid Dynamics (CFD).
Physical data storage is accomplished via a general database manager,
either \HDF (Hierarchical Data Format) or ADF (Advanced Data Format).
The present document describes the mapping from the Standard Interface
Data Structures (SIDS) to the file structure provided by the HDF
database manager.
A similar document, the \textit{SIDS-to-ADF File Mapping Manual},
describes the mapping for the ADF database manager.


\clearemptydoublepage
\section{CGNS Documentation}
\label{s:documentation}
\thispagestyle{plain}

Documentation of CGNS is found in a number of related
publications. These are maintained separately for several
reasons. First, they describe logically independent aspects of the CGNS
system. Second, many users will find that a subset of the documentation
is sufficient for their needs. And last, some portions of the system can
be viewed as independent entities useful outside the context of CGNS.

The main documents currently available are:

\begin{itemize*}
\item CGNS Overview and Entry-Level Document
\item A User's Guide to CGNS
\item CGNS Mid-Level Library
\item Standard Interface Data Structures
\item SIDS File Mapping Manual (this document)
\item CGIO User's Guide
\item ADF Implementation
\item HDF5 Implementation
\end{itemize*}

\subsection{Description of the Documents}

In order to make the current document as self-contained as possible,
basic information regarding all components of the CGNS standard has
been included. However, the reader should be aware of the documents
describing these other components. These should be regarded as the
authoritative and complete descriptions of their respective components
of CGNS. Each of these documents is described briefly in this section.

\subsubsection{CGNS Overview and Entry-Level Document}

This is an introductory document which provides an overall view of the
purpose and components of CGNS. The Overview is intended as an entry
point into CGNS. Prospective users completely unfamiliar with CGNS
should consult the Overview to determine whether CGNS provides the
capabilities they seek.

\subsubsection{A User's Guide to CGNS}

The \textit{User's Guide to CGNS} has been written to aid users in the
implementation of CGNS.
It is intended as a tutorial: light in content, but heavy in examples,
advice, and guidelines.
The guide provides a concise overview of many of the most commonly-used
features of the SIDS, and gives coding examples using the CGNS Mid-Level
Library to write and read simple CGNS files.

\subsubsection{The CGNS Mid-Level Library}

The \textit{CGNS Mid-Level Library} document describes a set of routines
which store and retrieve the CFD data objects defined in the SIDS.
Their purpose is to provide CGNS compliant I/O without the need for
detailed programming in the ADF or HDF Core.
These ``mid-level routines'' are designed to be inserted directly into
applications codes, such as flow solvers and grid generators.

\subsubsection{Standard Interface Data Structures}

The \textit{Standard Interface Data Structures}, usually abbreviated as
``SIDS,'' define the ``intellectual content'' of CFD data.
They detail the data which must be stored to completely characterize
each CFD entity. They describe, for example, what exactly is meant by a
``grid'' or a ``boundary condition''.
They also establish a system of nomenclature which gives standard
meaning to certain names, such as ``Density'' and ``SubsonicInflow''.

The SIDS description of the CFD data is hierarchical in nature, in that
complex entities are built up out of simpler ones. In the SIDS document,
this is reflected in a syntax which uses C-like structures to define
the various entities. The result is a tree-like structure which
exactly parallels the File Mapping described by the current
document. Thus in terms of
basic structure, the Mapping itself summarizes the SIDS. However, there
are many conventions regarding the nomenclature and meaning of data
which are not summarized in the current document, and for these the SIDS
is the authoritative document.

It is worth emphasizing that the SIDS may be regarded as a stand-alone
definition of the data associated with CFD, and that these data could
be stored in any sufficiently general format, given a mapping onto that
format.

\subsubsection{The SIDS File Mapping Manual}

The \textit{SIDS File Mapping Manual} specifies the exact manner
in which, under CGNS conventions, CFD data structures (the SIDS) are
to be stored in, i.e., mapped onto, the file structure provided by a
database manager.
Adherence to the mapping conventions guarantees uniform meaning
and location of CFD data within database files, and thereby allows the
construction of universal software to read and write the data.

\subsubsection{CGIO User's Guide}

The \textit{CGIO User's Guide} describes the interface between the
SIDS File Mapping and an underlying database manager.
It also documents the low-level routines that provide this
functionality as used in the Mid-Level Library.
This is the preferred way to access the database directly if
additional requirements beyond those provided by the MLL are needed. 

\subsubsection{ADF Implementation}

The \textit{ADF Implementation} describes one of the underlying database
managers, ADF (Advanced Data Format), which may be used to create CGNS
files.
ADF is a binary format, based on a simple tree structure.
In principal, nearly any kind of data could be stored in ADF format.
ADF was, however, especially designed for the storage of large
quantities of scientific data in a platform-independent and
randomly-accessible manner.
The ``ADF Core'' is a set of portable routines which store and retrieve
data in ADF format.
These routines are written in C, but Fortran versions are also provided.
The \textit{ADF User's Guide} describes both the ADF format and the ADF
Core routines in detail.

It should be emphasized that ADF, with its Core routines, constitutes a
very general stand-alone database manager which is not directly related
to CFD. It can therefore be used to store any kind of data \emph{once it has
been specified where to place that data within the ADF format}.

\subsubsection{HDF Implementation}

HDF5 (Hierarchical Data Format) is another underlying database manager
that may be used to create CGNS files, and is developed and maintained
by the National Center for Supercomputing Applications at the University
of Illinois at Urbana-Champaign.
Documentation for HDF5 is available at the HDF5 Home Page, at
\url{http://hdf.ncsa.uiuc.edu/HDF5/}.

Like ADF, HDF5 is a binary format, based on a tree structure,
designed for the storage of large quantities of scientific data in a
platform-independent and randomly-accessible manner.
Also like ADF, HDF5 can be used to store any kind of data.
The \textit{HDF5 Implementation} document describes the CGNS
conventions governing the placement of CFD-specific data within the HDF5
format.

\subsection{Which Documents Do You Need?}

Ideally, all users of the CGNS system will want to have all the
documents available for reference. However, many will find it possible
to begin to use the system effectively without reading all the documents
beforehand.  In fact, since the CGNS system is intended to minimize
interaction with underlying data structures, some users will find
they need very little knowledge of the system's internal workings. We
distinguish four classes of users who may wish to consult the CGNS
Documentation.

\subsubsection{Prospective Users}

Prospective users are presumably unfamiliar with CGNS. They will
probably wish to begin with the Overview, or, if they require more
detailed information, one or more of the various papers that have been
written describing CGNS.
Beyond that, most will find a quick read of this file mapping document
enlightening as to the logical form of the contents of CGNS files.
Browsing the figures in
\hyperref[s:figures]{Appendix~\ref*{s:figures}}
of this document, as well as the SIDS itself, will provide some feel
for the scope of the system.
The \textit{User's Guide to CGNS}, and the \textit{CGNS Mid-Level
Library} document should give an indication of what might be required to
implement CGNS in a given application.
Prospective users should probably not concern themselves with the
details of ADF or HDF.

\subsubsection{End Users}

The end user is the practitioner of CFD who generates the grids,
runs the flow codes and/or analyzes the results. For this user, CGNS
provides a mechanism for accumulating the output of the various
processes related to CFD, e.g., grid generation and flow solution,
and for making this output available to subsequent processes or for
archiving final results. For this user, a scan of the Overview will
sufficiently explain the overall workings of the system. This includes
end user responsibilities for matters not governed by CGNS, such as the
maintenance of files and directories. The end user will also find useful
the \textit{User's Guide to CGNS}, as well as
those portions of the SIDS which deal with standard nomenclature. AIAA
98-3007 may also be useful if more details about the capabilities of
CGNS are desired.

The end user is, by definition, not involved in the building of
CGNS-compliant applications code.

\subsubsection{Applications Code Developers}

The applications code developer builds or maintains code to support the
various sub-processes encountered in CFD, e.g., grid generation, flow
solution, post-processing, or flow visualization. The code developer
must be able to install CGNS compliant I/O. The most convenient method
for doing so is to utilize the CGNS Mid-Level Library.
The \textit{User's Guide to CGNS} is the starting point for learning to
use the Mid-Level Library to create and use CGNS files.
The \textit{CGNS Mid-Level Library} document itself should
also be considered essential.
This library of routines will perform the most common I/O operations in
a CGNS-compliant manner.
However, even when the CGNS Library suffices to implement all necessary
I/O, an understanding of the SIDS and the SIDS File Mapping documents
will be useful.
It will likely be necessary to consult the SIDS to determine the
precise meaning of the nomenclature.
In the unlikely case that the MLL does not provide all the
functionality needed by the application, the \textit{CGIO User's Guide}
may be consulted for information on the low-level access to the
CGNS database. 

\subsubsection{CGNS System Developers}

CGNS System development can be kept somewhat compartmentalized.
Developers responsible for the maintenance or building of supplements
to the ADF or HDF Core, need not concern themselves with documentation
other than the \textit{ADF User Guide} or the HDF5 documentation.
System developers wishing to add to the CGNS Mid-Level Library will need
all the documents.
Theoretical developments, such as extensions to the SIDS, may possibly
be undertaken with a knowledge of the SIDS alone, but such contributions
must also be added to the \textit{SIDS File Mapping}
before they can be implemented.

\subsection{How to Use This Document}

Those wishing to do more than simply browse this document will
find that the detailed information begins with the General File
Mapping Concepts in \autoref{s:general}. The detailed textual node
descriptions in \autoref{s:nodes} are more useful as reference than
as sequential literature. The best overall technical view of the
layout of CGNS files can be acquired by reference to the figures in
\hyperref[s:figures]{Appendix~\ref*{s:figures}}.


\clearemptydoublepage
\section{CGNS Background}
\label{s:background}
\thispagestyle{plain}

The information in this section is supplied for the sake of
completeness. It is identical with the information found in the
Overview.

\subsection{Purpose}

The purpose of CGNS is to provide a standard for recording and
recovering computer data associated with the numerical solution of the
equations of fluid dynamics. The format implemented by this standard is
(1) general, (2) portable, (3) expandable, and (4) durable.

The CGNS system consists of a collection of conventions, and software
implementing those conventions, for the storage and retrieval of
CFD (computational fluid dynamics) data. The system consists of two
parts: (1) a standard format for recording the data, and (2) software
that reads, writes and modifies data in that format.  The format is a
conceptual entity established by the documentation; the software is a
physical product supplied to enable developers to access and produce
data recorded in that format. The CGNS standard, applied through the use
of the supplied software, is intended to do the following:

\begin{itemize*}
\item facilitate the exchange of CFD data
      \begin{itemize*}
      \item between sites.
      \item between applications codes.
      \item across computing platforms.
      \end{itemize*}
\item stabilize the archiving of CFD data.
\end{itemize*}

\subsection{Participation and Brief History}

The CGNS project originated around 1994--1995 through a series of
meetings between Boeing and NASA that addressed improved means for
transferring NASA technology to industrial use. It was held that
a principal impediment to technology transfer was the disparity
in I/O formats employed by various flow codes, grid generators,
and so forth. The CGNS system was conceived as a means to promote
``plug-and-play'' CFD.

Agreement was reached to develop CGNS at Boeing, under NASA Contract
NAS1-20267, with active participation by a team of CFD researchers from

\begin{itemize*}
\item NASA Langley Research Center
\item NASA Glenn Research Center
\item NASA Ames Research Center
\item Boeing St-Louis (McDonnell-Douglas Corporation).
\item Boeing Commercial Airplane Group Aerodynamics
\item Boeing Commercial Airplane Group Propulsion
\item ICEM CFD Engineering Corporation of Berkeley, California
\end{itemize*}
Also participating in the discussions at various times have been
researchers from
\begin{itemize*}
\item Defense and Space Group and Environmental Systems
\item Arnold Engineering Development Center, representing the NPARC Alliance
\item Wright-Patterson Air Force Base
\end{itemize*}

\subsection{Scope}

The principal target of CGNS is the data normally associated with
compressible viscous flow (i.e., the Navier-Stokes equations), but the
standard is also applicable to subclasses such as Euler and potential
flows.

CGNS Version 1.0, released 5/15/98, was limited to problems described
by multiblock structured grids. Version 1.1 addresses grids,
flowfields, boundary conditions, and block-to-block connection
information. Also included are a number of auxiliary items,
including nondimensionalization, reference state, and equation set
specifications. The extension to time-dependent flows and unstructured
grids is addressed in Version 2. Also included are links between CGNS
data and CAD geometry. Any mix of the following types of field data can
be recorded:

\begin{itemize*}
\item nodal.
\item edge-centered.
\item face-centered.
\item cell-centered.
\end{itemize*}
Block connections can be of the following types:
\begin{itemize*}
\item contiguous (one-to-one).
\item abutting (patched mismatched).
\item overlapping (Chimera).
\end{itemize*}

Much of the standard and the software is applicable to computational
field physics in general. Disciplines other than fluid dynamics would
need to augment the data definitions and storage conventions, but the
fundamental database software, which provides platform independence, is
not specific to fluid dynamics.


\clearemptydoublepage
%
% last modif: 2004-12-08
%
% NEW CONCEPT: the ADF node is no more a real node, we say now CGNS node,
% because ADF itself exist at the same time with HDF ! 
%
% ****************************************************************************
% ** WE SHOULD NOT DOCUMENT THE MAPPING ASSUMING THE INTERFACE LIBRARY IS ADFH
% ** ANY HDF TOOL WOULD BE ABLE TO PRODUCE AN HDF5 STRUCTURE UNDERSTOOD AS A
% ** CGNS COMPLIANT STORAGE STRUCTURE 
% ****************************************************************************
%
% I don't know if we have to add some implementation details that would have
% an actual impact on interface. For example, the ways we create datasets 
% and dataspaces, always with default properties. Certainly we should locate
% the places where the properties could/would/should be set.
%
% See ADFH.c comments, such as:
%%   /*
%%    * The ADF documentation allows the dimension values to be changed
%%    * without affecting the data, so long as the data type and number
%%    * of dimensions are the same. With HDF5, we could emulate that by
%%    * using extendable data spaces (with chunking). However this only
%%    * allows the data size to increase, not decrease, and coming up
%%    * with a good value for chunking is difficult. Since changing the
%%    * dimension values without rewiting the data is not a common
%%    * operation, I decided to use fixed sizes, then buffer the data
%%    * in these rare cases.
%%    */
%
% Clearly, this is an implementation choice, but it has an impact on
% interface. The question is : ``Can somebody else implement an interface
% AND an HDF5 file compliant without the emulation described above ?''
%

\section{SLL Mapping to \HDF}
\label{s:hdfsummary}
\thispagestyle{plain}

% Cheating: almost replace ADF with HDF ;)

The purpose of the current document is to describe the way in which
CFD data is to be represented in an \HDF data tree. To do this, it is necessary
to first describe the \HDF data structure itself in some
detail. Therefore, a conceptual summary of \HDF is given here in order
to make the current document relatively independent, and to allow
the reader to focus on those aspects of \HDF which are essential to
understanding the file mapping.

Any \HDF file with a conformant mapping is \emph{CGNS/\HDF} compliant. The
mapping has been made using a \emph{per-node} basis. Instead of having a new
mapping dedicated to \HDF, we have made a mapping from the ADF nodes
to a set of \HDF groups. While this is not an optimal mapping, the use of 
such an \HDF node allows us to re-use the ADF API without change.

The \HDF documentation should be used as the authoritative
references to resolve any issues not covered by this summary.

\subsection{General Description of \HDF}

% Here, LOW level is a bit scornfull, we may change this term
The \HDF library provides CGNS with a low level storage system.  
This library is developed and maintained by the National Center
for Supercomputing Applications at the University of Illinois at
Urbana-Champaign.
It offers a
free storage system, with support for a very large number of platforms
used in the scientific world.  This system is in charge of all
physical features such as storage, compression, virtual access to data
and many other services also useful for the CGNS users.

\HDF allows a fine tuning of the physical devices, however, the \SLL
use of \HDF forces all \emph{property lists} to \texttt{DEFAULT}
values.

Refer to the HDF5 web site at \url{http://hdf.ncsa.uiuc.edu/HDF5} for
further information about \HDF.

\subsection{Dedicated \HDF Structures}
\label{s:adfnodestructure}

The \emph{CGNS node} is an \HDF group. It contains attributes, an optional
dataset and optional child groups. This structure is required and can be
extended by other implementors as long as the following requirements are
fullfilled.

In a CGNS tree, most of the nodes are \emph{normal} nodes. These nodes
can optionally contain data and have other nodes as children. 

There are special nodes, such as the root node and
the nodes managing the \emph{links}. These special link structures use
the \HDF mount system. It is not necessary for either the MLL 
or \SLL users to understand the \HDF mount system, but it should be 
understood by \SLL implementors.

\subsubsection{The CGNS Node Mapping}

The CGNS node uses \emph{groups} and \emph{attributes} elements of \HDF. 
A basic CGNS node is composed of elements listed in \autoref{t:groupnode}.
The attribute names have been enclosed in square brackets, because the
SLL uses private attributes with a leading \textit{space} character.
These brackets are not part of the names.

\begin{table}[htbp]
\centering
\caption[Basic Node]{\textbf{Basic Node}}
\label{t:groupnode}
\begin{tabular}{l >{\quad}l >{\quad}l}
\\ \hline\hline \\*[-2ex]
\bold{GROUP} &  &  
\\*[1ex] \hline\hline \\*[-2ex]
% is there a reason why types are not native ints/chars ?
REQUIRED ATTRIBUTE  & \fort{[name]}   & \fort{U8LE[33]} \\
REQUIRED ATTRIBUTE  & \fort{[label]}  & \fort{U8LE[33]} \\
REQUIRED ATTRIBUTE  & \fort{[type]}   & \fort{U8LE[3]} \\
REQUIRED ATTRIBUTE  & \fort{[ order]} & \fort{H5T\_NATIVE\_INT} \\
OPTIONAL DATASET    & \fort{[ data]}  & 
\\*[1ex] \hline\hline
\end{tabular}
\end{table}

All the children groups of a node are understood as \SLL children nodes of this
node, unless the group name starts with a \emph{space} character.
The \emph{dataset} contains the actual data of a node. In the case of
a node without data (i.e. \texttt{MT} type), the dataset does not exist.

In addition to the attributes of a \emph{normal} node, special attribute 
names are reserved for the root node (see 
\autoref{t:rootnode}). 
The name of the root node is \texttt{[HDF5 MotherNode]}, its label is
\texttt{[Root Node of HDF5 File]}, type is \texttt{MT}.

\begin{table}[htbp]
\centering
\caption[Root Node]{\textbf{Root Node}}
\label{t:rootnode}
\begin{tabular}{l >{\quad}l >{\quad}l}
\\ \hline\hline \\*[-2ex]
\bold{GROUP} & \fort{[/]} &  
\\*[1ex] \hline\hline \\*[-2ex]
% is there a reason why types are not native ints/chars ?
REQUIRED ATTRIBUTE  & \fort{[ version]} & \fort{U8LE[33]} \\
REQUIRED ATTRIBUTE  & \fort{[ format]}  & \fort{U8LE[33]}
\\*[1ex] \hline\hline
\end{tabular}
\end{table}

The node hosting the link entries is a special node named \texttt{[ mount]}
without attributes nor data. This group should be a child of the root node,
its absolute path is \texttt{[/ mount]}. The group \texttt{[ mount]} contains
one group for each mounted file refered to by a link. 
Each entry has the attributes described in \autoref{t:linknode}. 
The link node itself is a group with the special name \texttt{[ link]}. 
This implementation is discussed further in \autoref{s:linkmechs}.
% clear enough ?

\begin{table}[htbp]
\centering
\caption[Link Node]{\textbf{Link Node}}
\label{t:linknode}
\begin{tabular}{l >{\quad}l >{\quad}l}
\\ \hline\hline \\*[-2ex]
\bold{GROUP} & \fort{[/ mount/ *]} &  
\\*[1ex] \hline\hline \\*[-2ex]
REQUIRED ATTRIBUTE  & \fort{[ refcnt]} & \fort{H5T\_NATIVE\_INT} \\
REQUIRED ATTRIBUTE  & \fort{[ file]}   & \fort{H5T\_NATIVE\_CHAR}
\\*[1ex] \hline\hline
\end{tabular}
\end{table}

\subsubsection{The Node ID}
The node ID is a unique identifier assigned to each existing node by
\HDF when the file containing it is opened, and to new nodes as they are
created. \ADFH inquiries generally return node IDs as a result and
accept node IDs as input. By building a table of IDs, calling software
can subsequently access specific nodes without further search. The Node
ID is real and is not under user control, its lifetime is ended by
the closing of the \HDF file.

\subsubsection{The Node Name}
The node Name is a 32-byte character field which is user
controllable. Its general use is to distinguish among the children of a
given node; consequently, no two children of the same parent may have
the same Name. Constraints related to the node name are detailed in
\autoref{s:nodename}.

\subsubsection{The Label}
The Label is a 32-byte character field which is user controllable. \ADFH
assigns no formal role to the Label, but the intent was to identify
the structure of the included data. It is common for the various
children of a single parent to store different instances of the same
structure. Therefore, there is no prohibition against more than one
child of the same parent having the same Label.

\subsubsection{The Data Type}
The Data Type is a 32-byte character field which specifies the type and
precision of any data which is stored in the data field. Types provided
by \HDF are listed in \autoref{t:datatype}.

\begin{table}[htbp]
\centering
\caption[Data Types]{\textbf{Data Types}}
\label{t:datatype}
\begin{tabular}{l >{\quad}l >{\quad}l}
\\ \hline\hline \\*[-2ex]
\bold{Data Type} & \bold{Notation} & \bold{\HDF Type}
\\*[1ex] \hline\hline \\*[-2ex]
No Data             & \fort{MT} & \fort{-} \\
Integer 32          & \fort{I4} & \fort{H5T\_NATIVE\_INT32} \\
Integer 64          & \fort{I8} & \fort{H5T\_NATIVE\_INT64} \\
Unsigned Integer 32 & \fort{U4} & \fort{H5T\_NATIVE\_UINT32} \\
Unsigned Integer 64 & \fort{U8} & \fort{H5T\_NATIVE\_UINT64} \\
Real 32             & \fort{R4} & \fort{H5T\_NATIVE\_FLOAT} \\
Real 64             & \fort{R8} & \fort{H5T\_NATIVE\_DOUBLE} \\
Complex 64          & \fort{X4} & \fort{-} \\
Complex 128         & \fort{X8} & \fort{-} \\
Character           & \fort{C1} & \fort{H5T\_NATIVE\_CHAR} \\
Byte                & \fort{B1} & \fort{H5T\_NATIVE\_UCHAR} \\
Link                & \fort{LK} & \fort{-} 
\\*[1ex] \hline\hline
\end{tabular}
\end{table}

\noindent
%If the data type is \fort{MT} or \fort{LK}, the node attributes which
%described the data may be left undefined. 
There is no mapping to \HDF
\fort{MT} and \fort{LK} types, because there is no actual
\emph{data space} associated with the nodes. The type information itself
is stored in the node as the strings ``\fort{MT}'' and ``\fort{LK}''.
The types \fort{X4} and \fort{X8}
are also not mapped.

The data storage format is translated as described in \autoref{t:nativeformat}.
\begin{table}[htbp]
\centering
\caption[Native Formats]{\textbf{Native Formats}}
\label{t:nativeformat}
\begin{tabular}{l >{\quad}l}
\\ \hline\hline \\*[-2ex]
\bold{Native format} & \bold{\HDF Type}
\\*[1ex] \hline\hline \\*[-2ex]
\fort{H5T\_IEEE\_F32BE} & \fort{IEEE\_BIG\_32} \\
\fort{H5T\_IEEE\_F32LE} & \fort{IEEE\_LITTLE\_32} \\
\fort{H5T\_IEEE\_F64BE} & \fort{IEEE\_BIG\_64} \\
\fort{H5T\_IEEE\_F64LE} & \fort{IEEE\_LITTLE\_64}
\\*[1ex] \hline\hline
\end{tabular}
\end{table}

\subsubsection{The Number of Dimensions}

The Data portion of a node is designed to store multi-dimensional
arrays of data, each element of which is presumed to be of the Data
Type specified. The Number of Dimensions specifies the number of
integers required to reference a single datum within the array.

\subsubsection{The Dimension Values}

The Dimension Values are a list of integers expressing the actual
sizes of the stored array in each of the dimensions specified. These
dimensions are stored by the \emph{dataspace} associated
with the \emph{dataset}; no other attributes contains these values.

\subsubsection{The Data}

In an \HDF node, the portion of the node holding the actual stored data array
is a \textit{dataset}.

\subsubsection{The Child Table}

All groups of the current group are said to be the children of this group,
except the groups with a name starting with a space character. Each of
these children groups is a \emph{normal} node (i.e. group).

Children may be identified by their names and labels, and, thence,
by their node IDs once these have been determined. \HDF provides no
notion of order among children, but the \SLL layer adds a creation order
stored in the \texttt{[ order]} attribute. This order is guaranteed to
be the same from call to call, even after the file has been closed and
re-opened.

%In particular, the order of a list of
%children returned by ADF has nothing to do with the order in which they
%were inserted in the file. However, the order returned is consistent
%from call to call provided the file has not been closed and the node
%structure has not been modifed.

Note that there is no \emph{parent} table; that is, a node has no direct
knowledge of its parent. Since calling software must open the file
from the root, it presumably cannot access a child without having
first accessed the parent.  It is the responsibility of the calling
software to record the node ID of the parent if this information will be
required.

\subsubsection{The link mechanism}
\label{s:linkmechs}

A \texttt{LK} typed node is a \emph{link}. Such a node refers to
another node elsewhere. In other words, the \emph{link} has no child
or contents, but is a name of a node somewhere in the current file or
in another file.  The \texttt{LK} typed node is said to be the
\emph{emph} node, while the node elsewhere is said to be the
\emph{destination} node.

Is is the role of the \SLL layer to insure consistency\footnote{In
particular in order to avoid an acyclic graph.} and
transparency of the link mechanism, so that any \emph{normal} node
request to the \emph{link} node is performed as if it is performed
on the destination node.

A \emph{link} destination can be in the same file as the \emph{link} source
or in another file. In both cases, the \emph{link} is made using an \HDF
soft link from source to destination. In the case of a destination in another
file, the destination file is mounted in the \texttt{[/ mount]}\footnote{Absolute path.} group and
the soft link is made from the source to the destination now present in
this mounted file. 
% Ok. Is there still somebody in the room ?
% No ? Now, going to explain how to cook the pancakes.

Each time a file is mounted, the \texttt{[/ refcnt]}\footnote{Absolute path.}
attribute is incremented.
The file is unmounted if there is no reference to itself.




\clearemptydoublepage
\section{General Remarks}
\label{s:general}
\thispagestyle{plain}

\subsection{Acquiring the Software and Documentation}

The CGNS Mid-Level Library may be downloaded from SourceForge, at
{\itshape\url{http://sourceforge.net/projects/cgns}}.
This manual, as well as the other CGNS documentation, is available in
both HTML and PDF format from the CGNS documentation web site, at
{\itshape\url{http://www.grc.nasa.gov/www/cgns/CGNS_docs_current/}}.

\subsection{Organization of This Manual}

The sections that follow describe the Mid-Level Library functions
in detail.
The first three sections cover some basic file operations (i.e.,
opening and closing a CGNS file, and some configuration options)
(\autoref{s:fileops}), accessing a specific node in a CGNS database
(\autoref{s:navigating}), and error handling (\autoref{s:error}).
The remaining sections describe the functions used to read, write, and
modify nodes and data in a CGNS database.
These sections basically follow the organization used in
the ``Detailed CGNS Node Descriptions'' section of the
\href{../filemap/filemap.pdf}{SIDS File Mapping} manual.

At the start of each sub-section is a \textit{Node} line, listing the
the applicable CGNS node label.

Next is a table illustrating the syntax for the Mid-Level Library
functions.
The C functions are shown in the top half of the table, followed by the
corresponding Fortran routines in the bottom half of the table.
Input variables are shown in an \textcolor{input}{\texttt{upright blue}}
font, and output variables are shown in a
\textcolor{output}{\texttt{\textit{slanted red}}} font.
As of Version 3.1, some of the arguments to the Mid-Level Library have
changed from \texttt{int} to \texttt{cgsize\_t} in order to support 64-bit data.
For each function, the right-hand column lists the modes (read, write,
and/or modify) applicable to that function.

The input and output variables are then listed and defined.

\subsection{Language}

The CGNS Mid-Level Library is written in C, but each function has a
Fortran counterpart.
All function names start with ``\texttt{cg\_}''.
The Fortran functions have the same name as their C counterpart with the
addition of the suffix ``\texttt{\_f}''.

\subsection{Character Strings}

All data structure names and labels in CGNS are limited to 32
characters.
When reading a file, it is advised to pre-allocate the
character string variables to 32 characters in Fortran, and 33
in C (to include the string terminator).
Other character strings, such as the CGNS file name or
descriptor text, are unlimited in length.
The space for unlimited length character strings will be created by the
Mid-Level Library; it is then the responsibility of the application to
release this space by a call to \texttt{cg\_free}, described in
\autoref{s:free}.

\subsection{Error Status}

All C functions return an integer value representing the error
status.
All Fortran functions have an additional parameter, \texttt{ier},
which contains the value of the error status.
An error status different from zero implies that an error occured.
The error message can be printed using the error
handling functions of the CGNS library, described in \autoref{s:error}.
The error codes are coded in the C and Fortran include files
\textit{cgnslib.h} and \textit{cgnslib\_f.h}.

\subsection{Typedefs}
\label{s:typedefs}

Beginning with version 3.1, two new typedef variables have been
introduced to support 64-bit mode. The \texttt{cglong\_t} typedef
is always a 64-bit integer, and \texttt{cgsize\_t} will be either
a 32-bit or 64-bit integer depending on how the library was built.
Many of the C functions in the MLL have been changed
to to use \texttt{cgsize\_t} instead of \texttt{int} in the arguments.
These functions include any that may exceed the \texttt{2Gb} limit
of an \texttt{int}, e.g. zone dimensions, element data, boundary
conditions, and connectivity. In Fortran, all integer data is taken
to be \texttt{integer*4} for 32-bit and \texttt{integer*8} for 64-bit
builds.

Several types of variables are defined using typedefs in the
\textit{cgnslib.h} file.
These are intended to facilitate the implementation of
CGNS in C.
These variable types are defined as an enumeration of key
words admissible for any variable of these types.
The file \textit{cgnslib.h} must be included in any C application
programs which use these data types.

In Fortran, the same key words are defined as integer parameters in the
include file \textit{cgnslib\_f.h}.
Such variables should be declared as \texttt{integer} in Fortran
applications.
The file \textit{cgnslib\_f.h} must be included in any
Fortran application using these key words.

Note that the first two enumerated values in these lists,
\textit{xxxNull} and \textit{xxxUserDefined}, are only available
in the C interface, and are provided in the advent that your
C compiler does strict type checking. In Fortran, these values
are replaced by the numerically equivalent \texttt{CG\_Null}
and \texttt{CG\_UserDefined}. These values are also defined in the
C interface, thus either form may be used. The function prototypes
for the MLL use \texttt{CG\_Null} and \texttt{CG\_UserDefined},
rather than the more specific values.

The list of enumerated values (key words) for each of these variable
types (typedefs) are:
        
{\raggedright
\begin{Ventryi}{\texttt{AverageInterfaceType\_t}}
   \item [\texttt{ZoneType\_t}]
         \texttt{ZoneTypeNull, ZoneTypeUserDefined, Structured, Unstructured}
   \item [\texttt{ElementType\_t}]
         \texttt{ElementTypeNull, ElementTypeUserDefined, NODE, BAR\_2, BAR\_3, TRI\_3, TRI\_6, QUAD\_4,
         QUAD\_8, QUAD\_9, TETRA\_4, TETRA\_10, PYRA\_5,
         PYRA\_14, PENTA\_6, PENTA\_15, PENTA\_18, HEXA\_8,
         HEXA\_20, HEXA\_27, MIXED, PYRA\_13, NGON\_n, NFACE\_n}
   \item [\texttt{DataType\_t}]
         \texttt{DataTypeNull, DataTypeUserDefined, Integer, RealSingle, RealDouble, Character}
   \item [\texttt{DataClass\_t}]
         \texttt{DataClassNull, DataClassUserDefined, Dimensional, NormalizedByDimensional,
         NormalizedByUnknownDimensional, NondimensionalParameter,
         DimensionlessConstant}
   \item [\texttt{MassUnits\_t}]
         \texttt{MassUnitsNull, MassUnitsNullUserDefined, Kilogram, Gram, Slug, PoundMass}
   \item [\texttt{LengthUnits\_t}]
         \texttt{LengthUnitsNull, LengthUnitsUserDefined, Meter, Centimeter, Millimeter, Foot,
         Inch}
   \item [\texttt{TimeUnits\_t}]
         \texttt{TimeUnitsNull, TimeUnitsUserDefined, Second}
   \item [\texttt{TemperatureUnits\_t}]
         \texttt{TemperatureUnitsNull, TemperatureUnitsUserDefined, Kelvin, Celsius, Rankine,
         Fahrenheit}
   \item [\texttt{AngleUnits\_t}]
         \texttt{AngleUnitsNull, AngleUnitsUserDefined, Degree, Radian}
   \item [\texttt{ElectricCurrentUnits\_t}]
         \texttt{ElectricCurrentUnitsNull, ElectricCurrentUnitsUserDefined, Ampere, Abampere, Statampere, Edison,
         auCurrent}
   \item [\texttt{SubstanceAmountUnits\_t}]
         \texttt{SubstanceAmountUnitsNull, SubstanceAmountUnitsUserDefined, Mole, Entities, StandardCubicFoot,
         StandardCubicMeter}
   \item [\texttt{LuminousIntensityUnits\_t}]
         \texttt{LuminousIntensityUnitsNull, LuminousIntensityUnitsUserDefined, Candela, Candle, Carcel, Hefner, Violle}
   \item [\texttt{GoverningEquationsType\_t}]
         \texttt{GoverningEquationsTypeNull, GoverningEquationsTypeUserDefined, FullPotential, Euler, NSLaminar,
         NSTurbulent, NSLaminarIncompressible,
         NSTurbulentIncompressible}
   \item [\texttt{ModelType\_t}]
         \texttt{ModelTypeNull, ModelTypeUserDefined, Ideal, VanderWaals, Constant,
         PowerLaw, SutherlandLaw, ConstantPrandtl, EddyViscosity,
         ReynoldsStress, ReynoldsStressAlgebraic,
         Algebraic\_BaldwinLomax, Algebraic\_CebeciSmith,
         HalfEquation\_JohnsonKing, OneEquation\_BaldwinBarth,
         OneEquation\_SpalartAllmaras, TwoEquation\_JonesLaunder,
         TwoEquation\_MenterSST, TwoEquation\_Wilcox,
	 CaloricallyPerfect, ThermallyPerfect,
	 ConstantDensity, RedlichKwong,
	 Frozen, ThermalEquilib, ThermalNonequilib,
	 ChemicalEquilibCurveFit, ChemicalEquilibMinimization,
	 ChemicalNonequilib, EMElectricField, EMMagneticField,
         EMConductivity, Voltage, Interpolated,
         Equilibrium\_LinRessler, Chemistry\_LinRessler}
   \item [\texttt{GridLocation\_t}]
         \texttt{GridLocationNull, GridLocationUserDefined, Vertex, IFaceCenter, CellCenter, JFaceCenter,
         FaceCenter, KFaceCenter, EdgeCenter}
   \item [\texttt{GridConnectivityType\_t}]
         \texttt{GridConnectivityTypeNull, GridConnectivityTypeUserDefined, Overset, Abutting, Abutting1to1}
   \item [\texttt{PointSetType\_t}]
         \texttt{PointSetTypeNull, PointSetTypeUserDefined, PointList, PointRange, PointListDonor, PointRangeDonor,
         ElementList, ElementRange, CellListDonor}
   \item [\texttt{BCType\_t}]
         \texttt{BCTypeNull, BCTypeUserDefined, BCAxisymmetricWedge,
         BCDegenerateLine, BCExtrapolate,
         BCDegeneratePoint, BCDirichlet, BCFarfield, BCNeumann,
         BCGeneral, BCInflow, BCOutflow, BCInflowSubsonic,
         BCOutflowSubsonic, BCInflowSupersonic,
         BCOutflowSupersonic, BCSymmetryPlane, BCTunnelInflow,
         BCSymmetryPolar, BCTunnelOutflow, BCWallViscous,
         BCWall, BCWallViscousHeatFlux, BCWallInviscid,
         BCWallViscousIsothermal, FamilySpecified}
   \item [\texttt{BCDataType\_t}]
         \texttt{BCDataTypeNull, BCDataTypeUserDefined, Dirichlet, Neumann}
   \item [\texttt{RigidGridMotionType\_t}]
         \texttt{RigidGridMotionTypeNull, RigidGridMotionTypeUserDefined, ConstantRate, VariableRate}
   \item [\texttt{ArbitraryGridMotionType\_t}]
         \texttt{ArbitraryGridMotionTypeNull, ArbitraryGridMotionTypeUserDefined, NonDeformingGrid, DeformingGrid}
   \item [\texttt{SimulationType\_t}]
         \texttt{SimulationTypeNull, SimulationTypeUserDefined, TimeAccurate, NonTimeAccurate}
   \item [\texttt{WallFunctionType\_t}]
         \texttt{WallFunctionTypeNull, WallFunctionTypeUserDefined, Generic}
   \item [\texttt{AreaType\_t}]
         \texttt{AreaTypeNull, AreaTypeUserDefined, BleedArea, CaptureArea}
   \item [\texttt{AverageInterfaceType\_t}]
         \texttt{AverageInterfaceTypeNull, AverageInterfaceTypeUserDefined, AverageAll, AverageCircumferential, AverageRadial, AverageI,
	 AverageJ, AverageK}
\end{Ventryi}
}

\subsection{Character Names For Typedefs}

The CGNS library defines character arrays which map the typedefs above
to character strings.
These are global arrays dimensioned to the size of each list of
typedefs.
To retrieve a character string representation of a typedef, use the
typedef value as an index to the appropiate character array.
For example, to retrieve the string ``\texttt{Meter}''
for the \texttt{LengthUnits\_t Meter} typedef, use
\texttt{LengthUnitsName[Meter]}.
Functions are available to retrieve these names without the need for
direct global data access.
These functions also do bounds checking on the input, and if out of
range, will return the string ``\texttt{<invalid>}''.
An additional benefit is that these will work from within a Windows DLL,
and are thus the recommended access technique.
The routines have the same name as the global data arrays, but with a
``\texttt{cg\_}'' prepended.
For the example above, use ``\texttt{cg\_LengthUnitsName(Meter)}''.

\begin{table}
\begin{center}
\begin{tabular}{l}
\hline
Typedef Name Access Functions \\
\hline
\textcolor{output}{const char *name} = cg\_MassUnitsName(\textcolor{input}{MassUnits\_t type}); \\
\textcolor{output}{const char *name} = cg\_LengthUnitsName(\textcolor{input}{LengthUnits\_t type}); \\
\textcolor{output}{const char *name} = cg\_TimeUnitsName(\textcolor{input}{TimeUnits\_t type}); \\
\textcolor{output}{const char *name} = cg\_TemperatureUnitsName(\textcolor{input}{TemperatureUnits\_t type}); \\
\textcolor{output}{const char *name} = cg\_ElectricCurrentUnitsName(\textcolor{input}{ElectricCurrentUnits\_t type}); \\
\textcolor{output}{const char *name} = cg\_SubstanceAmountUnitsName(\textcolor{input}{SubstanceAmountUnits\_t type}); \\
\textcolor{output}{const char *name} = cg\_LuminousIntensityUnitsName(\textcolor{input}{LuminousIntensityUnits\_t type}); \\
\textcolor{output}{const char *name} = cg\_DataClassName(\textcolor{input}{DataClass\_t type}); \\
\textcolor{output}{const char *name} = cg\_GridLocationName(\textcolor{input}{GridLocation\_t type}); \\
\textcolor{output}{const char *name} = cg\_BCDataTypeName(\textcolor{input}{BCDataType\_t type}); \\
\textcolor{output}{const char *name} = cg\_GridConnectivityTypeName(\textcolor{input}{GridConnectivityType\_t type}); \\
\textcolor{output}{const char *name} = cg\_PointSetTypeName(\textcolor{input}{PointSetType\_t type}); \\
\textcolor{output}{const char *name} = cg\_GoverningEquationsTypeName(\textcolor{input}{GoverningEquationsType\_t type}); \\
\textcolor{output}{const char *name} = cg\_ModelTypeName(\textcolor{input}{ModelType\_t type}); \\
\textcolor{output}{const char *name} = cg\_BCTypeName(\textcolor{input}{BCType\_t type}); \\
\textcolor{output}{const char *name} = cg\_DataTypeName(\textcolor{input}{DataType\_t type}); \\
\textcolor{output}{const char *name} = cg\_ElementTypeName(\textcolor{input}{ElementType\_t type}); \\
\textcolor{output}{const char *name} = cg\_ZoneTypeName(\textcolor{input}{ZoneType\_t type}); \\
\textcolor{output}{const char *name} = cg\_RigidGridMotionTypeName(\textcolor{input}{RigidGridMotionType\_t type}); \\
\textcolor{output}{const char *name} = cg\_ArbitraryGridMotionTypeName(\textcolor{input}{ArbitraryGridMotionType\_t type}); \\
\textcolor{output}{const char *name} = cg\_SimulationTypeName(\textcolor{input}{SimulationType\_t type}); \\
\textcolor{output}{const char *name} = cg\_WallFunctionTypeName(\textcolor{input}{WallFunctionType\_t type}); \\
\textcolor{output}{const char *name} = cg\_AreaTypeName(\textcolor{input}{AreaType\_t type}); \\
\textcolor{output}{const char *name} = cg\_AverageInterfaceTypeName(\textcolor{input}{AverageInterfaceType\_t type}); \\
\hline
\end{tabular}
\end{center}
\end{table}
 
\clearpage

\subsection{64-bit C Portability and Issues}

If you use the \texttt{cgsize\_t} data type in new code,
it will work in both 32 and 64-bit compilation modes.
In order to support CGNS versions prior to 3.1, you may also
want to add something like this to your code:

\noindent \texttt{\#if CGNS\_VERSION < 3100}  \\
\noindent \texttt{\#define cgsize\_t int}  \\
\noindent \texttt{\#endif}

Existing code that uses \texttt{int} will not work with a
CGNS 3.1 library compiled in 64-bit mode. You may want
to add something like this to your code:

\noindent \texttt{\#if CGNS\_VERSION >= 3100 \&\& CG\_BUILD\_64BIT} \\
\noindent \texttt{\#error does not work in 64 bit mode} \\
\noindent \texttt{\#endif}

\noindent or modify your code to use \texttt{cgsize\_t}.

\subsection{64-bit Fortran Portability and Issues}

All integer arguments in the Fortran interface are taken to be
\texttt{integer*4} in 32-bit mode and \texttt{integer*8}
in 64-bit mode. If you have used default or implicit integers in
your Fortran code, it should port to 64-bit mode in most cases by
simply turning on your compiler option that promotes implicit
integers to \texttt{integer*8}. If you have explicitly defined your
integers as \texttt{integer*4}, your code will not work in 64-bit mode.
In that case, you will either need to change them to \texttt{integer}
(recommended for portability) or \texttt{integer*8}.

A new integer parameter has been added to the \textit{cgnslib\_f.h}
header, \texttt{CG\_BUILD\_64BIT}, which will be set
to 1 in 64-bit mode and 0 otherwise. You may use this parameter
to check at run time if the CGNS library has been compiled in 64-bit
mode or not, as in:

\noindent \texttt{if (CG\_BUILD\_64BIT .ne. 0) then} \\
\indent \texttt{print \*,'will not work in 64-bit mode'} \\
\indent \texttt{stop} \\
\noindent \texttt{endif}

If you are using a CGNS library prior to version 3.1, this parameter will
not be defined and you will need to rely on your compiler initializing all
undefined values to 0 (not always the case) for this test to work.

If your compiler supports automatic promotion of integers, and you use
implicit integers, your code should port to 64-bit with the following exception.

If you use an \texttt{Integer} data type in any routine that takes
a data type specification, and an implicit integer for the data, the code
will fail when compiled in 64-bit mode with automatic integer promotion.
An example of this would be:

\noindent \texttt{integer dim} \\
\noindent \texttt{integer data(dim)} \\
\noindent \texttt{call cg\_array\_write\_f('array',Integer,1,dim,data)} \\

This is because the MLL interprets the \texttt{Integer} data type
as \texttt{integer*4} regardless of the compilation mode. The compiler,
however, has automatically promoted \texttt{data} to be
\texttt{integer*8}. What you will need to do to prevent this problem,
is to either explicitly define \texttt{data} as in:

\noindent \texttt{integer dim} \\
\noindent \texttt{integer*4 data(dim)} \\
\noindent \texttt{call cg\_array\_write\_f('array',Integer,1,dim,data)} \\

\noindent or

\noindent \texttt{integer dim} \\
\noindent \texttt{integer*8 data(dim)} \\
\noindent \texttt{call cg\_array\_write\_f('array',LongInteger,1,dim,data)} \\

\noindent or test on \texttt{CG\_BUILD\_64BIT} as in:

\noindent \texttt{integer dim} \\
\noindent \texttt{integer data(dim)} \\
\noindent \texttt{if (CG\_BUILD\_64BIT .eq. 0) then} \\
\indent  \texttt{call cg\_array\_write\_f('array',Integer,1,dim,data)} \\
\noindent \texttt{else} \\
\indent  \texttt{call cg\_array\_write\_f('array',LongInteger,1,dim,data)} \\
\noindent \texttt{endif}

The last 2 options will only work with CGNS Version 3.1, since
\texttt{LongInteger} and \texttt{CG\_BUILD\_64BIT} are
not defined in previous versions.

You may also need to be careful when using integer constants as arguments
in 64-bit mode. If your compiler automatically promotes integer constants
to \texttt{integer*8}, then there is no problem. This is probably
the case if your compiler supports implicit integer promotion. If not,
then the constants will be \texttt{integer*4}, and your code
will not work in 64-bit mode. In that case you will need to do something like:

\noindent \texttt{integer*8 one,dim} \\
\noindent \texttt{integer*4 data(dim)} \\
\noindent \texttt{one = 1} \\
\noindent \texttt{call cg\_array\_write\_f('array',Integer,one,dim,data)}


\clearemptydoublepage
\section{Detailed CGNS Node Descriptions}
\label{s:nodes}
\thispagestyle{plain}

This section, together with the figures in
\hyperref[s:figures]{Appendix~\ref*{s:figures}}, constitutes a complete
description of the CGNS database structure, together with detailed
descriptions of the contents of each attribute of each node. It is
intended to be suitable as a reference for anyone implementing CGNS
using the ADF Core, but should also be of interest to those wishing to
understand exactly where information is stored within a CGNS database.

Note that it is the advertised purpose of the CGNS Mid-Level Library to
store and retrieve information in conformity with the mapping herein
described. Therefore, anyone accessing a CGNS database through the
CGNS routines alone does not need a detailed understanding of the file
mapping per se. However, this document should still prove useful in
ascertaining the meaning of some of the arguments which must be supplied
to the Library routines. Further, it will be necessary to consult the
SIDS themselves to determine some of the naming conventions.

The node descriptions in this section are closely
coupled to and cross-referenced with the figures in
\hyperref[s:figures]{Appendix~\ref*{s:figures}}, which show all the
nodes defined in the SIDS that have child nodes.
In the current section, the ``\textbf{Children:}'' entry in the list of
\textit{Node Attributes} is a reference to the figure showing that node
with its children.

The nodal hierarchy of the CGNS database directly reflects that of the
SIDS. Certain sections of the SIDS, notably Sections 4, 5, and 12,
describe basic data structures which appear repeatedly as children
of nodes representing more complex structures. In order to simplify
the presentation and avoid the introduction of undefined terms, this
section has been divided into two parts. \autoref{s:basicnodes}
defines a number of basic types which recur often in the structure,
and \autoref{s:specializednodes} describes higher level nodes of more
specific function built from those in \autoref{s:basicnodes}.\footnote{For
convenience, we group the node descriptions according to the type of
data the nodes contain. These groupings only roughly correspond to the
chapters in the SIDS.}

\subsection{Basic CGNS Nodes}
\label{s:basicnodes}

In this section we describe CGNS nodes which hold fundamental
\emph{types} of information. Their structure and function, which are the
same everywhere, are described here. However, the \emph{meaning} of the
data they hold at any particular point in a CGNS database depends on the
context, i.e., the parent node. Therefore, where necessary, any special
context-dependent meaning is elaborated in the paragraph devoted to the
parent.

\subsubsection{Descriptor Group}

These are user-assigned nodes designed to further describe the user's
intent. Their data is meant for human perusal or other user-designated
purposes.

\paragraph{\texttt{Descriptor\_t}}

The purpose of a node of type \fort{Descriptor\_t} is to store textual
data. It is not intended to be read by any system software, except to
return the text for human examination.

The intent of the \fort{Descriptor\_t} node is to hold comment sufficient
to allow someone other than the originator to understand what the
node contains. This may consist of a problem description, reference
documents, personal notes, etc.

Any node may have zero to many \fort{Descriptor\_t} children, with
the uses differentiated by the Name field. At this time, there are \emph{no}
conventions for these names or for the form of the associated text. It
is expected that a standard set such as \fort{README}, \fort{TimeStamp},
etc. will evolve as a matter of practice.

\textit{\uline{Node Attributes}}
\begin{Ventryic}{\textbf{Dimension Values:}}
\item [\textbf{Name:}]
      User defined
\item [\textbf{Label:}]
      \fort{Descriptor\_t}
\item [\textbf{DataType:}]
      \fort{C1}
\item [\textbf{Dimension:}]
      1
\item [\textbf{Dimension Values:}]
      Length of string to be stored, including any carriage control or
      null bytes.
\item [\textbf{Data:}]
      String - Since line terminators can be stored within the data,
      the user could conceptually store an entire document in this
      area, read it into a program and then print it out. For example,
      an entire PostScript document describing the problem (and maybe
      results) could be stored in the Data field, read by a program and
      then sent to a printer.
\item [\textbf{Children:}]
      None
\item [\textbf{Cardinality:}]
      0,$N$
\end{Ventryic}

\paragraph{\texttt{Ordinal\_t}}

Because there is no notion of order among children, there is
occasionally a desire to order children in a way that survives from one
opening of a CGNS database to another. The current CGNS Library provides
means of doing this. However, another early method was to place the node
``number'' in a child of type \fort{Ordinal\_t}.

Like \fort{Descriptor\_t}, the \fort{Ordinal\_t} node is completely
under the control of the user, who takes full responsibility for its
content. Unlike \fort{Descriptor\_t}, CGNS conventions do not encourage
the use of \fort{Ordinal\_t}, as it usually encodes information which
is redundant with the name. It is not read or written by standard
CGNS software, and so there is no assurance that sibling nodes
will be differently, consecutively, or consistently numbered by
\fort{Ordinal\_t}. Clearly, if \fort{Ordinal\_t} must be used, no node
should have more than one \fort{Ordinal\_t} child, and no two siblings
should have \fort{Ordinal\_t} children containing the same data.

It is worth noting that, if consistent numbering is desired, one way of
achieving it is to make the desired integer either the name or part of
the name. In fact, if, for example, individual zones are left unnamed,
the default convention will provide names of \fort{Zone1}, \fort{Zone2},
etc. Alternatively, the character strings ``\fort{1}'', ``\fort{2}'',
\ldots, are legal names. ADF or CGNS software, of course, will return
these as strings. This may necessitate type conversion or parsing before
the names can be used as integer indices.

\textit{\uline{Node Attributes}}
\begin{Ventryic}{\textbf{Dimension Values:}}
\item [\textbf{Name:}]
      \fort{Ordinal}
\item [\textbf{Label:}]
      \fort{Ordinal\_t}
\item [\textbf{DataType:}]
      \fort{I4}
\item [\textbf{Dimension:}]
      1
\item [\textbf{Dimension Values:}]
      1
\item [\textbf{Data:}]
      The user-defined ordinal number (an integer).
\item [\textbf{Children:}]
      None
\item [\textbf{Cardinality:}]
      0,1
\end{Ventryic}

\subsubsection{Physical Data Group}

\paragraph{\texttt{DataClass\_t}}

A \fort{DataClass\_t} node specifies the dimensional nature of the
data in and below its parent. It overrides any \fort{DataClass\_t}
information higher up in the tree. There are six recognized string
values. It is necessary to consult the SIDS to determine the precise
meaning.

\textit{\uline{Node Attributes}}
\begin{Ventryic}{\textbf{Dimension Values:}}
\raggedright
\item [\textbf{Name:}]
      \fort{DataClass}
\item [\textbf{Label:}]
      \fort{DataClass\_t}
\item [\textbf{DataType:}]
      \fort{C1}
\item [\textbf{Dimension:}]
      1
\item [\textbf{Dimension Values:}]
      The length of the string
\item [\textbf{Data:}]
      One of: \fort{DataClassNull}, \fort{DataClassUserDefined},
      \fort{Dimensional}, \fort{NormalizedByDimensional},
      \fort{NormalizedByUnknownDimensional},
      \fort{NondimensionalParameter}, or \fort{DimensionlessConstant}
\item [\textbf{Children:}]
      None
\item [\textbf{Cardinality:}]
      0,1
\end{Ventryic}

\paragraph{\texttt{DimensionalUnits\_t}}

A \fort{DimensionalUnits\_t} node specifies dimensional units
which apply to data in and below its parent.  It overrides any
\fort{DimensionalUnits\_t} information higher up in the tree. There are
five strings to specify, corresponding, respectively, to units for mass,
length, time, temperature, and angular measure. The number of recognized
string values varies with the physical property.

Units for three additional types of data are specified in a child node,
\fort{AdditionalUnits\_t}.

\textit{\uline{Node Attributes}}
\begin{Ventryic}{\textbf{Dimension Values:}}
\item [\textbf{Name:}]
      \fort{DimensionalUnits}
\item [\textbf{Label:}]
      \fort{DimensionalUnits\_t}
\item [\textbf{DataType:}]
      \fort{C1}
\item [\textbf{Dimension:}]
      2
\item [\textbf{Dimension Values:}]
      (32,5)
\item [\textbf{Data:}]
      \begin{Ventryc}{For temperature, one of:}
      \raggedright
      \item [For mass, one of:]
            \fort{MassUnitsNull}, \fort{MassUnitsUserDefined},
            \fort{Kilogram}, \fort{Gram}, \fort{Slug}, \fort{Pound-Mass}
      \item [For length, one of:]
            \fort{LengthUnitsNull}, \fort{LengthUnitsUserDefined},
            \fort{Meter}, \fort{Centimeter}, \fort{Millimeter}, \fort{Foot},
            \fort{Inch}
      \item [For time, one of:]
            \fort{TimeUnitsNull}, \fort{TimeUnitsUserDefined}, \fort{Second}
      \item [For temperature, one of:]
            \fort{TemperatureUnitsNull}, \fort{TemperatureUnitsUserDefined},
            \fort{Kelvin}, \fort{Celsius}, \fort{Rankine}, \fort{Fahrenheit}
      \item [For angles, one of:]
            \fort{AngleUnitsNull}, \fort{AngleUnitsUserDefined},
            \fort{Degree}, \fort{Radian}
      \end{Ventryc}
\item [\textbf{Children:}]
      See \autoref{f:DimensionalUnits}
\item [\textbf{Cardinality:}]
      0,1
\end{Ventryic}

\paragraph{\texttt{AdditionalUnits\_t}}

An \fort{AdditionalUnits\_t} node specifies dimensional units for
additional types of data.
To maintain compatibility with earlier CGNS versions, this is an
optional child node of \fort{DimensionalUnits\_t}.
The specified units apply to data in and below the parent of the
corresponding \fort{DimensionalUnits\_t} node, and override any
\fort{AdditionalUnits\_t} information higher up in the tree.
There are three strings to specify, corresponding, respectively, to
units for electric current, substance amount, and luminous intensity.
The number of recognized string values varies with the physical
property.

\textit{\uline{Node Attributes}}
\begin{Ventryic}{\textbf{Dimension Values:}}
\item [\textbf{Name:}]
      \fort{AdditionalUnits}
\item [\textbf{Label:}]
      \fort{AdditionalUnits\_t}
\item [\textbf{DataType:}]
      \fort{C1}
\item [\textbf{Dimension:}]
      2
\item [\textbf{Dimension Values:}]
      (32,3)
\item [\textbf{Data:}]
      \begin{Ventryc}{For luminous intensity, one of:}
      \raggedright
      \item [For electric current, one of:]
            \fort{ElectricCurrentUnitsNull},
            \fort{ElectricCurrentUnitsUserDefined}, \fort{Ampere},
            \fort{Abampere}, \fort{Statampere}, \fort{Edison},
            \fort{auCurrent}
      \item [For substance amount, one of:]
            \fort{SubstanceAmountUnitsNull},
            \fort{SubstanceAmountUnitsUserDefined}, \fort{Mole},
            \fort{Entities}, \fort{StandardCubicFoot},
            \fort{StandardCubicMeter}
      \item [For luminous intensity, one of:]
            \fort{LuminousIntensityUnitsNull},
            \fort{LuminousIntensityUnitsUserDefined}, \fort{Candela},
            \fort{Candle}, \fort{Carcel}, \fort{Hefner},
            \fort{Violle}
      \end{Ventryc}
\item [\textbf{Children:}]
      None
\item [\textbf{Cardinality:}]
      0,1
\end{Ventryic}

\paragraph{\texttt{DataConversion\_t}}

A \fort{DataConversion\_t} node specifies a non-homogeneous linear
function which converts non-dimensional data in its parent to
raw dimensional data. Although in principle it overrides any
\fort{DataConversion\_t} information higher up in the tree, it is
generally not meaningful for it to apply to more than one kind of
physical data. Therefore, CGNS specifies its use only as a child of a
node which actually contains a single type of real data.

There are two values to specify, corresponding to the scale factor and
offset. The SIDS contain the exact conversion formula.

\textit{\uline{Node Attributes}}
\begin{Ventryic}{\textbf{Dimension Values:}}
\item [\textbf{Name:}]
      \fort{DataConversion}
\item [\textbf{Label:}]
      \fort{DataConversion\_t}
\item [\textbf{DataType:}]
      \fort{R4} or \fort{R8}
\item [\textbf{Dimension:}]
      1
\item [\textbf{Dimension Values:}]
      2
\item [\textbf{Data:}]
      \fort{ConversionScale}, \fort{ConversionOffset}
\item [\textbf{Children:}]
      None
\item [\textbf{Cardinality:}]
      0,1
\end{Ventryic}

\paragraph{\texttt{DimensionalExponents\_t}}

A \fort{DimensionalExponents\_t} node specifies the powers of mass,
length, time, temperature, and angular measure which characterize
dimensional data in its parent. Although in principle it overrides any
\fort{DimensionalExponents\_t} information higher up in the tree, it
is generally not meaningful for it to apply to more than one kind of
physical data. Therefore, CGNS specifies its use only as a child of
a node which actually contains a single type of real data.
There are five values to specify, corresponding to the five types of
units specified using \fort{DimensionalUnits\_t}.
The data type is real, not integer.

Exponents for three additional types of data are specified in a child node,
\fort{AdditionalExponents\_t}.

\textit{\uline{Node Attributes}}
\begin{Ventryic}{\textbf{Dimension Values:}}
\raggedright
\item [\textbf{Name:}]
      \fort{DimensionalExponents}
\item [\textbf{Label:}]
      \fort{DimensionalExponents\_t}
\item [\textbf{DataType:}]
      \fort{R4} or \fort{R8}
\item [\textbf{Dimension:}]
      1
\item [\textbf{Dimension Values:}]
      5
\item [\textbf{Data:}]
      \fort{MassExponent}, \fort{LengthExponent}, \fort{TimeExponent},
      \fort{TemperatureExponent}, \fort{AngleExponent}
\item [\textbf{Children:}]
      See \autoref{f:DimensionalExponents}
\item [\textbf{Cardinality:}]
      0,1
\end{Ventryic}

\paragraph{\texttt{AdditionalExponents\_t}}

An \fort{AdditionalExponents\_t} node specifies the powers of the units
for additional types of data, which characterize the corresponding
dimensional data.
There are three values to specify, corresponding to the three types of
units specified using \fort{AdditionalUnits\_t}.
The data type is real, not integer.

\textit{\uline{Node Attributes}}
\begin{Ventryic}{\textbf{Dimension Values:}}
\raggedright
\item [\textbf{Name:}]
      \fort{AdditionalExponents}
\item [\textbf{Label:}]
      \fort{AdditionalExponents\_t}
\item [\textbf{DataType:}]
      \fort{R4} or \fort{R8}
\item [\textbf{Dimension:}]
      1
\item [\textbf{Dimension Values:}]
      3
\item [\textbf{Data:}]
      \fort{ElectricCurrentExponent}, \fort{SubstanceAmountExponent},
      \fort{LuminousIntensityExponent}
\item [\textbf{Children:}]
      None
\item [\textbf{Cardinality:}]
      0,1
\end{Ventryic}

\paragraph{\texttt{DataArray\_t}}

A \fort{DataArray\_t} node is a very general type of node meant to hold
large arrays of data, such as grids and flowfields. Often, some of the
attributes of a \fort{DataArray\_t} node depend on the context in which
the node is found; that is, they are structure parameters.

For example, the SIDS specify that the Data Type of \fort{DataArray\_t}
is a structure parameter, ``\fort{DataType}'', which may assume
any of the values ``\fort{I\textit{n}}'', ``\fort{R\textit{n}}'',
or ``\fort{C\textit{n}}''.

The other two attributes of \fort{DataArray\_t}, \fort{Dimensions}
and \fort{DataSize}, also depend on the context where they are being used.
\fort{Dimensions} is a function of the underlying dimensionality of
the data being described (often \fort{IndexDimension}, defined in the
\fort{CGNSBase\_t} node), and the \fort{DataSize} may be inferred from
detailed descriptions of the grid.

A node may have any number of \fort{DataArray\_t} children. The
meaning of their contents is differentiated by Name, often
according to conventions specified by the SIDS. SIDS names are
usually precise and descriptive, such as \fort{CoordinateTheta} or
\fort{EnergyInternal}. (For a current list of sanctioned names, see the
SIDS, Appendix A.) Conversely, quantities not specified by the SIDS can
be stored in \fort{DataArray\_t} nodes, but should be given names other
than those specified in the SIDS. In other words, to comply with the
SIDS requires that one give a quantity the SIDS-defined name \emph{if
and only if} it is one of the SIDS-defined quantities.

\textit{\uline{Node Attributes}}
\begin{Ventryic}{\textbf{Dimension Values:}}
\item [\textbf{Name:}]
      Context dependent
\item [\textbf{Label:}]
      \fort{DataArray\_t}
\item [\textbf{DataType:}]
      Context dependent
\item [\textbf{Dimension:}]
      Context dependent
\item [\textbf{Dimension Values:}]
      Context dependent
\item [\textbf{Data:}]
      The array of data values
\item [\textbf{Children:}]
      See \autoref{f:DataArray}
\item [\textbf{Cardinality:}]
      0,$N$
\item [\textbf{Parameters:}]
      \fort{DataType}, dimension of the data, size of the data
\end{Ventryic}

\paragraph{Integer Arrays}

Integer array nodes perform the same function as nodes of type
\fort{DataArray\_t}, but store integer instead of real arrays. They are
always of type \fort{int[]} or \ital{cgsize\_t}\fort{[]},
with the dimensions and values given either
explicitly in the appropriate fields, or as parameters or functions.
The \ital{cgsize\_t} data type is used wherever
there is the potential of exceeding the size limit for a
32-bit integer (\fort{I4}), such as \fort{PointList},
\fort{PointRange}, and \fort{ElementRange}.

\textit{\uline{Node Attributes}}
\begin{Ventryic}{\textbf{Dimension Values:}}
\raggedright
\item [\textbf{Name:}]
      Context dependent
\item [\textbf{Label:}]
      \fort{int}, \fort{int[IndexDimension]}, \fort{int[2*IndexDimension]},
      or \fort{int[1 + ... + IndexDimension]}
\item [\textbf{DataType:}]
      \fort{I4} or \ital{cgsize\_t}
\item [\textbf{Dimension:}]
      1
\item [\textbf{Dimension Values:}]
      1, \fort{IndexDimension}, \fort{2*IndexDimension}, or
      \fort{(1 + ... + IndexDimension)}
\item [\textbf{Data:}]
      The array of integer values
\item [\textbf{Cardinality:}]
      0,1
\item [\textbf{Children:}]
      None
\item [\textbf{Parameters:}]
      \fort{IndexDimension} or none (context dependent)
\end{Ventryic}

\subsubsection{Location and Position Group}

\paragraph{\texttt{GridLocation\_t}}

A \fort{GridLocation\_t} node specifies the physical location, with
respect to the underlying grid, with which the field data below
its parent is associated. The value (data field) is a character
string of enumeration type, i.e., it must take one of a number of
predefined values. These values are: \fort{Vertex}, \fort{CellCenter},
\fort{FaceCenter}, \fort{IFaceCenter}, \fort{JFaceCenter},
\fort{KFaceCenter}, or \fort{EdgeCenter}. The strings are case
sensitive, and an exact match is required. The \fort{GridLocation\_t}
node is optional, and the default is \fort{Vertex}.

\newpage
\textit{\uline{Node Attributes}}
\begin{Ventryic}{\textbf{Dimension Values:}}
\raggedright
\item [\textbf{Name:}]
      \fort{GridLocation}
\item [\textbf{Label:}]
      \fort{GridLocation\_t}
\item [\textbf{DataType:}]
      \fort{C1}
\item [\textbf{Dimension:}]
      1
\item [\textbf{Dimension Values:}]
      Length of the string value
\item [\textbf{Data:}]
      \fort{Vertex}, \fort{CellCenter}, \fort{FaceCenter}, \fort{IFaceCenter},
      \fort{JFaceCenter}, \fort{KFaceCenter}, or \fort{EdgeCenter}
\item [\textbf{Children:}]
      None
\item [\textbf{Cardinality:}]
      0,1
\end{Ventryic}

\paragraph{\texttt{Rind\_t}}

The presence of a \fort{Rind\_t} node indicates that field data stored
below its parent includes values associated with a spatial extent beyond
that of the basic underlying grid. Such data often arise from the use of
ghost cells, or from the copying of information from adjacent zones.

Within a single zone, the size of the basic grid is found in the data
field of the \fort{Zone\_t} node (see \autoref{s:Zone}).
The data field of a \fort{Rind\_t} node contains integers specifying
the number of planes (for structured grids) or number of rind points or
elements (for unstructured grids) of included extra data.
The planes for structured grids correspond to the low and high values in
the $i$-direction, low and high values in the $j$-direction, and low and
high values in the $k$-direction (if needed), in that order.
Note that the actual size of the field data, which is stored in a
\fort{DataArray\_t} sibling node, is a \fort{DataSize} structure
parameter which depends on the basic grid size, the \fort{GridLocation},
and the \fort{Rind}.

The \fort{Rind\_t} node is optional, and the default is no rind.

\textit{\uline{Node Attributes}}
\begin{Ventryic}{\textbf{Dimension Values:}}
\item [\textbf{Name:}]
      \fort{Rind}
\item [\textbf{Label:}]
      \fort{Rind\_t}
\item [\textbf{DataType:}]
      \fort{I4}
\item [\textbf{Dimension:}]
      1
\item [\textbf{Dimension Values:}]
      \fort{2*IndexDimension}
\item [\textbf{Data:}]
      Number of planes of extra data in low $i$, high $i$, low $j$, high
      $j$, etc. (for structured grids) or number of points or elements
      of extra data (for unstructured grids)
\item [\textbf{Children:}]
      None
\item [\textbf{Cardinality:}]
      0,1
\item [\textbf{Parameters:}]
      \fort{IndexDimension}
\end{Ventryic}

\paragraph{\texttt{IndexRange\_t}}

An \fort{IndexRange\_t} node describes a subregion of a zone. This may
be, for example, a sub-block or a portion of a face of a zone. It may be
used to describe the locations of boundary condition patches and holes
for overset grids.

\textit{\uline{Node Attributes}}
\begin{Ventryic}{\textbf{Dimension Values:}}
\item [\textbf{Name:}]
      \fort{PointRange}, \fort{PointRangeDonor}, \fort{ElementRange},
      or user defined
\item [\textbf{Label:}]
      \fort{IndexRange\_t}
\item [\textbf{DataType:}]
      \ital{cgsize\_t}
\item [\textbf{Dimension:}]
      2
\item [\textbf{Dimension Values:}]
      \fort{IndexDimension}, 2
\item [\textbf{Data:}]
      First indices, last indices
\item [\textbf{Children:}]
      None
\item [\textbf{Cardinality:}]
      Context dependent
\item [\textbf{Parameters:}]
      \fort{IndexDimension}
\end{Ventryic}

\paragraph{\texttt{IndexArray\_t}}

An \fort{IndexArray\_t} node describes a general subregion of a
zone. Unlike \fort{In\-dex\-Range\_t}, it lists all the elements of the
subregion, rather then only the first and last ones. Its use is similar
to \fort{IndexRange\_t}.

\textit{\uline{Node Attributes}}
\begin{Ventryic}{\textbf{Dimension Values:}}
\raggedright
\item [\textbf{Name:}]
      \fort{PointList}, \fort{PointListDonor}, \fort{CellListDonor},
      or \fort{InwardNormalList}
\item [\textbf{Label:}]
      \fort{IndexArray\_t}
\item [\textbf{DataType:}]
      \ital{cgsize\_t}, or (for \fort{InwardNormalList}) \fort{R4} or \fort{R8}
\item [\textbf{Dimension:}]
      2
\item [\textbf{Dimension Values:}]
      \fort{IndexDimension}, number of items in the list;
      or (for \fort{InwardNormalList}) \fort{PhysicalDimension}, number
      of items in the list
\item [\textbf{Data:}]
      Index coordinates of each point or element in the list, or
      (for \fort{InwardNormalList}) physical-space normal vectors
      at each point or element in the list
\item [\textbf{Children:}]
      None
\item [\textbf{Cardinality:}]
      0,1
\item [\textbf{Parameters:}]
      \fort{IndexDimension}, either \fort{PointListSize} or
      \fort{ListLength}, and \fort{DataType};
      or (for \fort{InwardNormalList})
      \fort{PhysicalDimension}, \fort{ListLength}, and \fort{DataType}
\end{Ventryic}

\subsubsection{Auxiliary Data Group}

\paragraph{\texttt{ReferenceState\_t}}

The appearance of a \fort{ReferenceState\_t} node is optional. It is
used to specify the values of flow quantities at reference conditions,
e.g., at freestream or stagnation. This is typically done for the whole
database, in which case the \fort{ReferenceState\_t} node is a child of
the \fort{CGNSBase\_t} node.

\fort{ReferenceState\_t} nodes follow the usual convention that
information specified lower in the tree overrides higher level
specifications. Such overrides are therefore specified if a
\fort{ReferenceState\_t} node appears as a child of a \fort{Zone\_t},
\fort{ZoneBC\_t}, or \fort{BCDataSet\_t} node.

The actual values are stored in one or more \fort{DataArray\_t} children
whose names identify the quantities being stored. If present, the
units specified in the \fort{DimensionalUnits\_t} child apply to all
\fort{DataArray\_t} children, subject to the usual override convention.
( I.e., if one of the \fort{DataArray\_t} children itself has a
\fort{DimensionalUnits\_t} child, it takes precedence over the higher
level specification.)

\textit{\uline{Node Attributes}}
\begin{Ventryic}{\textbf{Dimension Values:}}
\item [\textbf{Name:}]
      \fort{ReferenceState}
\item [\textbf{Label:}]
      \fort{ReferenceState\_t}
\item [\textbf{DataType:}]
      \fort{MT}
\item [\textbf{Children:}]
      See \autoref{f:ReferenceState}
\item [\textbf{Cardinality:}]
      0,1
\end{Ventryic}

\paragraph{\texttt{ConvergenceHistory\_t}}

\fort{ConvergenceHistory\_t} nodes are intended for the storage of lists
of quantities accumulated during calculations associated with either the
entire CGNS database or with a single zone.

In the former case, they are called Global convergence histories, and
appear as children of the \fort{CGNSBase\_t} node. In the latter, they
are called Local and stored below, with the zones to which they correspond.

Each \fort{ConvergenceHistory\_t} node is a parent of a collection
of one-dimensional \fort{DataArray\_t} nodes, each of which contains
a list of values of a quantity defined by the user. These quantities
are differentiated by their user-assigned Names. User definitions of
the names are recorded in a \fort{Descriptor\_t} child node with Name
\fort{NormDefinitions}. Children of types \fort{DataClass\_t} and
\fort{DimensionalUnits\_t} modify the meaning of the \fort{DataArray\_t}
children in the usual manner.

\textit{\uline{Node Attributes}}
\begin{Ventryic}{\textbf{Dimension Values:}}
\raggedright
\item [\textbf{Name:}]
      \fort{GlobalConvergenceHistory} if under a \fort{CGNSBase\_t} node;
      \fort{ZoneConvergenceHistory} if under a \fort{Zone\_t} node
\item [\textbf{Label:}]
      \fort{ConvergenceHistory\_t}
\item [\textbf{DataType:}]
      \fort{I4}
\item [\textbf{Dimension:}]
      1
\item [\textbf{Dimension Values:}]
      1
\item [\textbf{Data:}]
      Number of iterations
\item [\textbf{Children:}]
      See \autoref{f:ConvergenceHistory}
\item [\textbf{Cardinality:}]
      0,1
\end{Ventryic}

\paragraph{\texttt{IntegralData\_t}}

\fort{IntegralData\_t} nodes are intended for the storage of
integrated flow quantities such as mass flows, forces and
moments. These are kept in \fort{DataArray\_t} children just as in the
\fort{ConvergenceHistory\_t} nodes, except that these nodes hold only
one real number each.

\textit{\uline{Node Attributes}}
\begin{Ventryic}{\textbf{Dimension Values:}}
\item [\textbf{Name:}]
      User defined
\item [\textbf{Label:}]
      \fort{IntegralData\_t}
\item [\textbf{DataType:}]
      \fort{MT}
\item [\textbf{Children:}]
      See \autoref{f:IntegralData}
\item [\textbf{Cardinality:}]
      0,$N$
\end{Ventryic}

\paragraph{\texttt{UserDefinedData\_t}}

\fort{UserDefinedData\_t} nodes are intended as a means of
storing arbitrary user-defined data in \fort{Descriptor\_t} and
\fort{DataArray\_t} children without the restrictions or implicit
meanings imposed on these node types at other node locations.

Multiple \fort{Descriptor\_t} and \fort{DataArray\_t} children
may be stored below a \fort{UserDefinedData\_t} node, and the
\fort{DataArray\_t} children may be of any dimension and size.

\textit{\uline{Node Attributes}}
\begin{Ventryic}{\textbf{Dimension Values:}}
\item [\textbf{Name:}]
      User defined
\item [\textbf{Label:}]
      \fort{UserDefinedData\_t}
\item [\textbf{DataType:}]
      \fort{MT}
\item [\textbf{Children:}]
      See \autoref{f:UserDefinedData}
\item [\textbf{Cardinality:}]
      0,$N$
\end{Ventryic}

\paragraph{\texttt{Gravity\_t}}

An optional \fort{Gravity\_t} node may be used to define the
gravitational vector.

\textit{\uline{Node Attributes}}
\begin{Ventryic}{\textbf{Dimension Values:}}
\item [\textbf{Name:}]
      \fort{Gravity}
\item [\textbf{Label:}]
      \fort{Gravity\_t}
\item [\textbf{DataType:}]
      \fort{MT}
\item [\textbf{Children:}]
      See \autoref{f:Gravity}
\item [\textbf{Cardinality:}]
      0,1
\end{Ventryic}

\subsection{Specialized Nodes}
\label{s:specializednodes}

In this section we describe nodes whose use is specialized to certain
types of CFD-related data. Although these nodes may appear in multiple
places in a CGNS DataBase, they play a single role in the description of
the data.

\subsubsection{Grid Specification}

CGNS recognizes the notion of a collection of subdomains called zones,
within each of which there is a single structured or unstructured
grid. Mathematically, the grid is an assignment of a location in
physical space to each element in a discrete computational space. An
essential feature of the grid is the connection structure it inherits
from the underlying computational space.

It is possible, given a grid, to create others from it, by translation
to cell centers, for example. However, CGNS views these as new
field structures associated with the original grid, and the File
Mapping specifies that they be stored as \fort{FlowSolution\_t} or
\fort{DiscreteData\_t} nodes (see \autoref{s:fieldspecification}).

\paragraph{\texttt{GridCoordinates\_t}}

A \fort{GridCoordinates\_t} node describes a grid associated with
a single zone. For a structured zone, the connection structure of
the underlying computational space is that of a rectangular array,
and its dimension is the \fort{IndexDimension}, that is, the number
of integers required to identify a point in the grid. The physical
dimension is the number of real coordinates assigned at each grid point
and need not be the same. Thus CGNS can store a grid, for example, with
\fort{IndexDimension} equal to two and a physical dimension of three,
that is, a structured grid on a curved surface.

\fort{IndexDimension} is a zone dependent parameter. For an unstructured
grid, it always equals one, meaning that a unique index is required to
specified a node location. For a structured grid, \fort{IndexDimension}
varries with the \fort{CellDimension} of the mesh. For a mesh composed
of 3D cells, \fort{IndexDimension} equals 3, while for a mesh composed
of surface or shell elements, \fort{IndexDimension} equals 2. The
values of the physical coordinates of the grid points are stored
in \fort{DataArray\_t} children of \fort{GridCoordinates\_t}. The
names of the coordinates are stored in the Name field of the
corresponding \fort{DataArray\_t} node. For common coordinate systems,
i.e., Cartesian, polar, cylindrical, and spherical, the names are
specified by the SIDS.

Unlike \fort{FlowSolution\_t} and \fort{DiscreteData\_t} nodes (see
\autoref{s:fieldspecification}), \fort{GridCoordinates\_t} nodes are
not permitted to have \fort{GridLocation\_t} children, because all grid
points are at vertices by definition.

Coordinate arrays may also contain rind data. If
they do, the \fort{GridCoordinates\_t} node must have a \fort{Rind\_t}
child node describing the amount of rind. All \fort{DataArray\_t} nodes
under \fort{GridCoordinates\_t} must have the same size. Because the
number of field quantities to be stored depends on the number of rind,
the actual dimension values are functions, specified in this document by
the generic term \fort{DataSize[]}.

Under each node of type \fort{Zone\_t}, the original grid is
contained in a node named \fort{GridCoordinates}.  Additional
\fort{GridCoordinates\_t} data structures are allowed, with user-defined
names, to store the grid at multiple time steps or iterations.

\textit{\uline{Node Attributes}}
\begin{Ventryic}{\textbf{Dimension Values:}}
\item [\textbf{Name:}]
      \fort{GridCoordinates} or user defined
\item [\textbf{Label:}]
      \fort{GridCoordinates\_t}
\item [\textbf{DataType:}]
      \fort{MT}
\item [\textbf{Children:}]
      See \autoref{f:GridCoordinates}
\item [\textbf{Cardinality:}]
      0,$N$
\item [\textbf{Parameters:}]
      \fort{IndexDimension}, \fort{VertexSize}
\item [\textbf{Functions:}]
      \fort{DataSize}
\end{Ventryic}

\paragraph{\texttt{Elements\_t}}

The \fort{Elements\_t} data structure is required for unstructured
zones, and contains the element connectivity data, the element type,
the element range, the parent elements data, and the number of boundary
elements.

\textit{\uline{Node Attributes}}
\begin{Ventryic}{\textbf{Dimension Values:}}
\item [\textbf{Name:}]
      User defined
\item [\textbf{Label:}]
      \fort{Elements\_t}
\item [\textbf{DataType:}]
      \fort{I4}
\item [\textbf{Dimension:}]
      1
\item [\textbf{Dimension Values:}]
      2
\item [\textbf{Data:}]
      \fort{ElementType} value, \fort{ElementSizeBoundary}
\item [\textbf{Children:}]
      See \autoref{f:Elements}
\item [\textbf{Cardinality:}]
      0,$N$
\end{Ventryic}

\paragraph{\texttt{Axisymmetry\_t}}

The \fort{Axisymmetry\_t} data structure may be included as a child
of the \fort{CGNSBase\_t} node to record the axis of rotation and the
angle of rotation around this axis for an axisymmetric database.

\textit{\uline{Node Attributes}}
\begin{Ventryic}{\textbf{Dimension Values:}}
\item [\textbf{Name:}]
      \fort{Axisymmetry}
\item [\textbf{Label:}]
      \fort{Axisymmetry\_t}
\item [\textbf{DataType:}]
      \fort{MT}
\item [\textbf{Children:}]
      See \autoref{f:Axisymmetry}
\item [\textbf{Cardinality:}]
      0,1
\end{Ventryic}

\paragraph{\texttt{RotatingCoordinates\_t}}

The \fort{RotatingCoordinates\_t} data structure may be included as a
child of either the \fort{CGNSBase\_t} node or a \fort{Zone\_t} node
to record the center of rotation and the rotation rate vector for a
rotating coordinate system.

\textit{\uline{Node Attributes}}
\begin{Ventryic}{\textbf{Dimension Values:}}
\item [\textbf{Name:}]
      \fort{RotatingCoordinates}
\item [\textbf{Label:}]
      \fort{RotatingCoordinates\_t}
\item [\textbf{DataType:}]
      \fort{MT}
\item [\textbf{Children:}]
      See \autoref{f:RotatingCoordinates}
\item [\textbf{Cardinality:}]
      0,1
\end{Ventryic}

\subsubsection{Field Specification}
\label{s:fieldspecification}

The object of computational field physics is to compute fields
of physical data associated with points in space.

\paragraph{\texttt{FlowSolution\_t}}

A \fort{FlowSolution\_t} node describes a field of physical data
associated with the grid for a single zone. It is intended for
the storage of computed flowfield data such as densities and
pressures. There is no convention as to how many or what kind
of quantities must or may be stored. In particular, it is not
specified that the quantities need in any sense be either complete or
non-redundant.

The data are stored in \fort{DataArray\_t} children of
\fort{FlowSolution\_t}. These \fort{DataArray\_t} nodes are dimensioned
by the same underlying \fort{IndexDimension} parameter as the grid, and
the order of storage within the \fort{DataArray\_t} nodes is presumed
the same as it is for the grid. The names of the physical quantities are
stored in the Name field of the corresponding \fort{DataArray\_t}
node. For common fluid dynamic quantities the names are specified by the
SIDS.

The relationship between the locations of the field quantities and the
vertices of the grid is specified by a \fort{GridLocation\_t} child
node. If this node is absent, the field quantities are assumed to be
associated with the grid vertices. Field arrays
may also contain rind data. If they do, the \fort{FlowSolution\_t} node
must have a \fort{Rind\_t} child node describing the amount of rind. All
\fort{DataArray\_t} nodes under a single \fort{FlowSolution\_t} must
have the same size. Field arrays containing different numbers of rind
must be stored under different \fort{FlowSolution\_t} nodes. There
may be any number of nodes of type \fort{FlowSolution\_t} under a
\fort{Zone\_t}.

Because the number of field quantities to be stored depends on the
number of rind and on the location with respect to the grid, the actual
dimension values are functions, specified in this document by the
generic term \fort{DataSize[]}.

The meaning of the field arrays is modified in the usual way by any
\fort{DataClass\_t} or \fort{DimensionalUnits\_t} children of the
\fort{FlowSolution\_t} node.

\textit{\uline{Node Attributes}}
\begin{Ventryic}{\textbf{Dimension Values:}}
\item [\textbf{Name:}]
      User defined
\item [\textbf{Label:}]
      \fort{FlowSolution\_t}
\item [\textbf{DataType:}]
      \fort{MT}
\item [\textbf{Children:}]
      See \autoref{f:FlowSolution}
\item [\textbf{Cardinality:}]
      0,$N$
\item [\textbf{Parameters:}]
      \fort{IndexDimension}, \fort{VertexSize}, \fort{CellSize}
\item [\textbf{Functions:}]
      \fort{DataSize}
\end{Ventryic}

\paragraph{\texttt{DiscreteData\_t}}

\fort{DiscreteData\_t} nodes are identical to \fort{FlowSolution\_t}
nodes, but are intended for the storage of fields of real data not
usually identified as part of the field solution, such as cell-centered
grids.

\textit{\uline{Node Attributes}}
\begin{Ventryic}{\textbf{Dimension Values:}}
\item [\textbf{Name:}]
      User defined
\item [\textbf{Label:}]
      \fort{DiscreteData\_t}
\item [\textbf{DataType:}]
      \fort{MT}
\item [\textbf{Children:}]
      See \autoref{f:FlowSolution}
\item [\textbf{Cardinality:}]
      0,$N$
\item [\textbf{Parameters:}]
      \fort{IndexDimension}, \fort{VertexSize}, \fort{CellSize}
\item [\textbf{Functions:}]
      \fort{DataSize}
\end{Ventryic}

\paragraph{\texttt{ZoneSubRegion\_t}}

The \fort{ZoneSubRegion\_t} node allows for the ability to give flowfield or
other information over a subset of the entire zone in a CGNS file.
This subset may be over a portion of a boundary, or it may be over a portion
of the entire field.

\textit{\uline{Node Attributes}}
\begin{Ventryic}{\textbf{Dimension Values:}}
\item [\textbf{Name:}]
      User defined
\item [\textbf{Label:}]
      \fort{ZoneSubRegion\_t}
\item [\textbf{DataType:}]
      \fort{I4}
\item [\textbf{Dimension:}]
      1
\item [\textbf{Dimension Values:}]
      \fort{RegionCellDimension}
\item [\textbf{Children:}]
      See \autoref{f:ZoneSubRegion}
\item [\textbf{Cardinality:}]
      0,$N$
\item [\textbf{Parameters:}]
      \fort{IndexDimension}, \fort{VertexSize}, \fort{CellSize}
\item [\textbf{Functions:}]
      \fort{DataSize}
\end{Ventryic}

\subsubsection{Connectivity Group}

\paragraph{\texttt{Transform} Node}

The \fort{Transform} node is a node of type \fort{int[]} which is
identified by its name rather than its label. Thus the name must be
``\fort{Transform}''. It appears only as a child of a node of type
\fort{GridConnectivity1to1\_t}.

This node stores the transformation matrix relating the indices of two
adjacent zones. Its data field contains a list of \fort{IndexDimension}
signed integers, each within the range \fort{\mbox{[-IndexDimension},
..., +IndexDimension]}, and no two of which have the same absolute
value. Thus in 3-D allowed components are 0, $\pm 1$, $\pm 2$, and
$\pm 3$. Each component of the array shows the image in the adjacent zone
of a positive index increment in the current zone. The SIDS contain
complete details.

\textit{\uline{Node Attributes}}
\begin{Ventryic}{\textbf{Dimension Values:}}
\item [\textbf{Name:}]
      \fort{Transform}
\item [\textbf{Label:}]
      ``\fort{int[IndexDimension]}''
\item [\textbf{DataType:}]
      \fort{I4}
\item [\textbf{Dimension:}]
      1
\item [\textbf{Dimension Values:}]
      \fort{IndexDimension}
\item [\textbf{Data:}]
      Transformation matrix (shorthand)
\item [\textbf{Children:}]
      None
\item [\textbf{Cardinality:}]
      0,1
\item [\textbf{Parameters:}]
      \fort{IndexDimension}
\end{Ventryic}

\paragraph{\texttt{GridConnectivityType\_t}}

The purpose of this node is to describe the type of zone-to-zone
connectivity specified by its parent, which is always a
\fort{GridConnectivity\_t} node. The connectivity type is given in the
data field as a character string which may take one of three specific
values: \fort{Abutting}, \fort{Abutting1to1}, or \fort{Overset}.

There is a shorthand form of the \fort{GridConnectivity\_t}
node, namely, \fort{GridConnectivity1to1\_t}, which
incorporates the assumption that the connection is
\fort{Abutting1to1}. Nodes of type \fort{GridConnectivity1to1\_t}
do not have \fort{GridConnectivityType\_t} subnodes. However,
\fort{GridConnectivity1to1\_t} nodes can only be used to specify
zone-to-zone connections on rectangular subregions between two
structured zones. So the use of \fort{GridConnectivityType\_t} subnodes
to specify \fort{Abutting1to1} is required if the connecting regions
are not rectangular, or if the connectivity involves a least one
unstructured zone.

\textit{\uline{Node Attributes}}
\begin{Ventryic}{\textbf{Dimension Values:}}
\item [\textbf{Name:}]
      \fort{GridConnectivityType}
\item [\textbf{Label:}]
      \fort{GridConnectivityType\_t}
\item [\textbf{DataType:}]
      \fort{C1}
\item [\textbf{Dimension:}]
      1
\item [\textbf{Dimension Values:}]
      Length of string
\item [\textbf{Data:}]
      \fort{Abutting}, \fort{Abutting1to1}, or \fort{Overset}
\item [\textbf{Children:}]
      None
\item [\textbf{Cardinality:}]
      0,1
\end{Ventryic}

\paragraph{\texttt{GridConnectivity1to1\_t}}

This node is a shorthand format of \fort{GridConnectivity\_t}
(see \autoref{s:GridConnectivity}) capable of describing only
\fort{Abutting1to1} connections between two structured zones. The
underlying subregion must have rectangular data structure.

Each \fort{GridConnectivity1to1\_t} node describes a subregion
of a face of a zone whose vertices are coincident in a 1-to-1
fashion with those of a corresponding subregion of a face of another
zone. Each \fort{ZoneGridConnectivity\_t} node may have as many
\fort{GridConnectivity1to1\_t} (or \fort{GridConnectivity\_t}) children
as are required to describe the connection structure.

The location of the connected subregion of a face of the current zone
is given in a single child of type \fort{IndexRange\_t}, whose name
is specified by the mapping as ``\fort{PointRange}''. The location of
the corresponding subregion on a face of the other zone is given in
a single child of type \fort{IndexRange\_t}, whose name is specified
by the mapping as ``\fort{PointRangeDonor}''. The first (i.e.,
beginning) points in these \fort{IndexRange\_t} nodes are presumed to be
coincident. The specification of the correspondence is completed by the
inclusion of a \fort{Transform} child node which describes the relative
orientation of the two systems of indices. The second (i.e., end) point
of the \fort{PointRange} subnode specifies the extant of the connection.

In general, the File Mapping seeks to avoid the storage of
redundant data. However, there are two redundancies associated with
\fort{GridConnectivity1to1\_t}. First, for the sake of symmetry, the
information recorded here is duplicated (in reverse) in a corresponding
node under the donor zone. It is expected that these two specifications
will agree.

Second, the end point of the \fort{PointRangeDonor} can be
calculated from the other three points specified, along with the
transform. However, the transform cannot be inferred from the four
points. Therefore, the end point of the \fort{PointRangeDonor} is
considered to be redundant, and the three points and the transform are
designated as the primary specification.

\textit{\uline{Node Attributes}}
\begin{Ventryic}{\textbf{Dimension Values:}}
\item [\textbf{Name:}]
      User defined
\item [\textbf{Label:}]
      \fort{GridConnectivity1to1\_t}
\item [\textbf{DataType:}]
      \fort{C1}
\item [\textbf{Dimension:}]
      1
\item [\textbf{Dimension Values:}]
      Length of string
\item [\textbf{Data:}]
      \fort{ZoneDonorName}
\item [\textbf{Children:}]
      See \autoref{f:GridConnectivity1to1}
\item [\textbf{Cardinality:}]
      0,$N$
\item [\textbf{Parameters:}]
      \fort{IndexDimension}
\end{Ventryic}

\paragraph{\texttt{GridConnectivity\_t}}
\label{s:GridConnectivity}

The \fort{GridConnectivity\_t} node is the most general format for
describing grid connectivity. It can describe one-to-one, mismatched,
and overset connectivity, and the underlying subregions of the
connecting zones need not be rectangular.

Each \fort{GridConnectivity\_t} node describes a subregion
of a zone which corresponds to a subregion of another
zone. Each \fort{ZoneGridConnectivity\_t} node may have as many
\fort{GridConnectivity\_t} (or \fort{GridConnectivity1to1\_t}) children
as are required to describe the connection structure.

The location of the connected subregion of the current zone is
given in a single child of type either \fort{IndexRange\_t} or
\fort{IndexArray\_t}, whose name is specified by the mapping as
``\fort{PointRange}'' or ``\fort{PointList}'', respectively.

If the grid connectivity is one-to-one, the corresponding subregion is
defined with a single child of type \fort{IndexArray\_t}, whose name is
specified by the mapping as ``\fort{PointListDonor}''.  Otherwise, the
corresponding subregion is defined by two child nodes, one defining the
cells and the other the interpolation factors within the cells.
See the SIDS for the complete description.

\textit{\uline{Node Attributes}}
\begin{Ventryic}{\textbf{Dimension Values:}}
\item [\textbf{Name:}]
      User defined
\item [\textbf{Label:}]
      \fort{GridConnectivity\_t}
\item [\textbf{DataType:}]
      \fort{C1}
\item [\textbf{Dimension:}]
      1
\item [\textbf{Dimension Values:}]
      Length of string
\item [\textbf{Data:}]
      \fort{ZoneDonorName}
\item [\textbf{Children:}]
      See \autoref{f:GridConnectivity}
\item [\textbf{Cardinality:}]
      0,$N$
\item [\textbf{Parameters:}]
      \fort{IndexDimension}, \fort{CellDimension}
\item [\textbf{Functions:}]
      \fort{PointListSize}
\end{Ventryic}

\paragraph{\texttt{GridConnectivityProperty\_t}}

An optional \fort{GridConnectivityProperty\_t} node may be used to
record special properties associated with particular connectivity
patches.

\textit{\uline{Node Attributes}}
\begin{Ventryic}{\textbf{Dimension Values:}}
\item [\textbf{Name:}]
      \fort{GridConnectivityProperty}
\item [\textbf{Label:}]
      \fort{GridConnectivityProperty\_t}
\item [\textbf{DataType:}]
      \fort{MT}
\item [\textbf{Children:}]
      See \autoref{f:GridConnectivityProperty}
\item [\textbf{Cardinality:}]
      0,1
\end{Ventryic}

\paragraph{\texttt{Periodic\_t}}

A \fort{Periodic\_t} node may be used as a child of
\fort{GridConnectivityProperty\_t} to record data associated with a
periodic interface.

\textit{\uline{Node Attributes}}
\begin{Ventryic}{\textbf{Dimension Values:}}
\item [\textbf{Name:}]
      \fort{Periodic}
\item [\textbf{Label:}]
      \fort{Periodic\_t}
\item [\textbf{DataType:}]
      \fort{MT}
\item [\textbf{Children:}]
      See \autoref{f:Periodic}
\item [\textbf{Cardinality:}]
      0,1
\end{Ventryic}

\paragraph{\texttt{AverageInterface\_t}}

An \fort{AverageInterface\_t} node is used as a child of
\fort{GridConnectivityProperty\_t} when data at the current connectivity
interface is to be averaged in some way prior to passing it to a
neighboring interface.

\textit{\uline{Node Attributes}}
\begin{Ventryic}{\textbf{Dimension Values:}}
\item [\textbf{Name:}]
      \fort{AverageInterface}
\item [\textbf{Label:}]
      \fort{AverageInterface\_t}
\item [\textbf{DataType:}]
      \fort{MT}
\item [\textbf{Children:}]
      See \autoref{f:AverageInterface}
\item [\textbf{Cardinality:}]
      0,1
\end{Ventryic}

\paragraph{\texttt{OversetHoles\_t}}

A node of type \fort{OversetHoles\_t} describes a region in a grid in
which solution values are to be ignored because the data in the region
is to be represented by values associated with other ``overlapping''
zones (equivalent to that specified by \fort{IBLANK} $= 0$ in the PLOT3D
format). Each \fort{ZoneGridConnectivity\_t} node may have as many
\fort{OversetHoles\_t} children as are required to describe the affected
region.

Each hole is described either by a single child of type \fort{IndexArray\_t} or
by any number of children of type \fort{IndexRange\_t}. The latter is provided
as a means of specifying holes which are unions of small numbers of
logically rectangular subregions. However, if the region is irregular,
the intent is that it should be specified by a single child of type
\fort{IndexArray\_t} which lists the points.

\textit{\uline{Node Attributes}}
\begin{Ventryic}{\textbf{Dimension Values:}}
\item [\textbf{Name:}]
      User defined
\item [\textbf{Label:}]
      \fort{OversetHoles\_t}
\item [\textbf{DataType:}]
      \fort{MT}
\item [\textbf{Children:}]
      See \autoref{f:OversetHoles}
\item [\textbf{Cardinality:}]
      0,$N$
\item [\textbf{Parameters:}]
      \fort{IndexDimension}
\end{Ventryic}

\paragraph{\texttt{ZoneGridConnectivity\_t}}

Each \fort{Zone\_t} node may have zero or more children of type
\fort{ZoneGridConnectivity\_t}. It holds no data, but serves as the point
below which all connectivity data associated with the zone can be found.
Multiple \fort{ZoneGridConnectivity\_t} nodes may be used to specify
time-dependent changes in the connectivity information.

\textit{\uline{Node Attributes}}
\begin{Ventryic}{\textbf{Dimension Values:}}
\item [\textbf{Name:}]
      \fort{ZoneGridConnectivity} or user defined
\item [\textbf{Label:}]
      \fort{ZoneGridConnectivity\_t}
\item [\textbf{DataType:}]
      \fort{MT}
\item [\textbf{Children:}]
      See \autoref{f:ZoneGridConnectivity}
\item [\textbf{Cardinality:}]
      0,$N$
\item [\textbf{Parameters:}]
      \fort{IndexDimension}, \fort{CellDimension}
\end{Ventryic}

\subsubsection{Boundary Condition Group}

Nodes in this group are used to specify the physical boundary
conditions. Each boundary condition is associated with a subregion of
a zone. For brevity below, we use the word ``domain'' to refer to the
region on which a boundary condition is to be enforced.

The domain is usually, but not necessarily, a subregion of a face of the
zone. The mapping is sufficiently general to permit the description of
internal boundary conditions and boundary conditions which do not lie on
a constant coordinate plane.

Mathematical boundary conditions are generally applied on subregions of
physical dimension one less than the corresponding field problem. This
condition, however, is neither defined nor enforced by the File Mapping.

A large number of standard boundary condition types are named by the
SIDS. In addition, it is possible to define new types as collections
of Dirichlet and Neumann conditions. It is not possible to describe
the entire array of possibilities within this document, and the reader
should consult the SIDS for a full description.

\paragraph{\texttt{InwardNormalIndex}}

An \fort{InwardNormalIndex} node is a node of type
\fort{int[IndexDimension]} which is identified by its Name.
It applies to structured grids only, and its function is to specify on
which side of the domain the condition is to be enforced.

\fort{InwardNormalIndex} may have only one nonzero element, whose
sign indicates the computational-coordinate direction of the boundary
condition patch normal; this normal points into the interior of the zone.
For example, if the domain lies on the face of a three-dimensional zone
where the second index is a maximum, the inward normal index values are
$[0,-1,0]$.

The \fort{InwardNormalIndex} node must apply to the entire domain
of the boundary condition.

For a boundary condition on a face of a zone, the
\fort{InwardNormalIndex} can be calculated from other data and need not
be specified. Its purpose is to define the normal direction for internal
boundary conditions and other cases where the direction is ambiguous.

\textit{\uline{Node Attributes}}
\begin{Ventryic}{\textbf{Dimension Values:}}
\item [\textbf{Name:}]
      \fort{InwardNormalIndex}
\item [\textbf{Label:}]
      ``\fort{int[IndexDimension]}''
\item [\textbf{DataType:}]
      \fort{I4}
\item [\textbf{Dimension:}]
      1
\item [\textbf{Dimension Values:}]
      \fort{IndexDimension}
\item [\textbf{Data:}]
      Index of inward normal
\item [\textbf{Children:}]
      None
\item [\textbf{Cardinality:}]
      0,1
\item [\textbf{Parameters:}]
      \fort{IndexDimension}
\end{Ventryic}

\paragraph{\texttt{InwardNormalList}}

An \fort{InwardNormalList} node is a node of type \fort{IndexArray\_t}
identified by its Name. Its data field contains an array of
physical (real) vectors which point into the region on which the
boundary condition is to be applied.  It may be used for boundary
conditions on complex domains for which \fort{InwardNormalIndex} is not
defined, or to store vectors orthogonal to the domain of the boundary
condition where these are not easily calculated from the domain itself.

\textit{\uline{Node Attributes}}
\begin{Ventryic}{\textbf{Dimension Values:}}
\item [\textbf{Name:}]
      \fort{InwardNormalList}
\item [\textbf{Label:}]
      \fort{IndexArray\_t}
\item [\textbf{DataType:}]
      \fort{R4} or \fort{R8}
\item [\textbf{Dimension:}]
      2
\item [\textbf{Dimension Values:}]
      \fort{PhysicalDimension}, \fort{ListLength}
\item [\textbf{Data:}]
      Inward normal vectors
\item [\textbf{Children:}]
      None
\item [\textbf{Cardinality:}]
      0,1
\item [\textbf{Parameters:}]
      \fort{PhysicalDimension}, \fort{ListLength}
\end{Ventryic}

\paragraph{\texttt{BCData\_t}}

When global or local Dirichlet or Neumann boundary conditions are
defined, a node of type \fort{BCData\_t} is introduced to store the
numerical data. For global data, this consists of a single quantity kept
in a \fort{DataArray\_t} child. For local
data, e.g., a pressure profile, it is a vector of quantities stored in
an order corresponding to that defining the domain and kept in a child
node of type \fort{DataArray\_t}.

\textit{\uline{Node Attributes}}
\begin{Ventryic}{\textbf{Dimension Values:}}
\item [\textbf{Name:}]
      \fort{DirichletData} or \fort{NeumannData}
\item [\textbf{Label:}]
      \fort{BCData\_t}
\item [\textbf{DataType:}]
      \fort{MT}
\item [\textbf{Children:}]
      See \autoref{f:BCData}
\item [\textbf{Cardinality:}]
      0,1
\item [\textbf{Parameters:}]
      \fort{ListLength}
\end{Ventryic}

\paragraph{\texttt{BCDataSet\_t}}

The function of a \fort{BCDataSet\_t} node is to specify the equations
to be applied at the boundary, including any actual data values which
may be required. The type of the equation is specified by the SIDS and
recorded in the data field. For some types, the data is implicit or
empty. For others, the data is specified in \fort{BCData\_t} children.

If the locations at which the boundary conditions are to be applied
are specified in \fort{BCDataSet\_t}, using \fort{PointRange} or
\fort{PointList}, the structure function \fort{ListLength} is used.
Otherwise, the structure parameter \fort{ListLength} is required.

\textit{\uline{Node Attributes}}
\begin{Ventryic}{\textbf{Dimension Values:}}
\item [\textbf{Name:}]
      User defined
\item [\textbf{Label:}]
      \fort{BCDataSet\_t}
\item [\textbf{DataType:}]
      \fort{C1}
\item [\textbf{Dimension:}]
      1
\item [\textbf{Dimension Values:}]
      Length of string
\item [\textbf{Data:}]
      \fort{BCTypeSimple} value
\item [\textbf{Children:}]
      See \autoref{f:BCDataSet}
\item [\textbf{Cardinality:}]
      0,$N$
\item [\textbf{Functions:}]
      \fort{ListLength}
\item [\textbf{Parameters:}]
      \fort{ListLength}
\end{Ventryic}

\paragraph{\texttt{BC\_t}}

A \fort{BC\_t} node specifies a single boundary condition to be applied
on a single zone. It specifies the domain on which the condition is
to be applied and the equations to be enforced. All the \fort{BC\_t}
nodes for a single zone are found under that zone's \fort{ZoneBC\_t}
node. A \fort{ZoneBC\_t} node may have as many \fort{BC\_t} children
as are required to describe the physical boundary conditions on the
corresponding zone.

The domain on which the boundary condition is to be enforced is
specified by a single node of type either \fort{IndexRange\_t} or
\fort{IndexArray\_t}. The equations are specified in one or
more \fort{BCDataSet\_t} children.

The type of the boundary condition, which may be either simple or
compound, is specified in the data field.  For a complete description,
it is necessary to consult the SIDS.

\textit{\uline{Node Attributes}}
\begin{Ventryic}{\textbf{Dimension Values:}}
\item [\textbf{Name:}]
      User defined
\item [\textbf{Label:}]
      \fort{BC\_t}
\item [\textbf{DataType:}]
      \fort{C1}
\item [\textbf{Dimension:}]
      1
\item [\textbf{Dimension Values:}]
      Length of string
\item [\textbf{Data:}]
      \fort{BCType} value
\item [\textbf{Children:}]
      See \autoref{f:BC}
\item [\textbf{Cardinality:}]
      0,$N$
\item [\textbf{Parameters:}]
      \fort{IndexDimension}, \fort{PhysicalDimension}
\end{Ventryic}

\paragraph{\texttt{ZoneBC\_t}}

The \fort{ZoneBC\_t} node occurs at most once for each zone and serves
as the location under which all boundary conditions on that zone are
collected.

\textit{\uline{Node Attributes}}
\begin{Ventryic}{\textbf{Dimension Values:}}
\item [\textbf{Name:}]
      \fort{ZoneBC}
\item [\textbf{Label:}]
      \fort{ZoneBC\_t}
\item [\textbf{DataType:}]
      \fort{MT}
\item [\textbf{Children:}]
      See \autoref{f:ZoneBC}
\item [\textbf{Cardinality:}]
      0,1
\item [\textbf{Parameters:}]
      \fort{IndexDimension}, \fort{PhysicalDimension}
\end{Ventryic}

\paragraph{\texttt{BCProperty\_t}}

An optional \fort{BCProperty\_t} node may be used to record special
properties associated with particular boundary condition patches.

\textit{\uline{Node Attributes}}
\begin{Ventryic}{\textbf{Dimension Values:}}
\item [\textbf{Name:}]
      \fort{BCProperty}
\item [\textbf{Label:}]
      \fort{BCProperty\_t}
\item [\textbf{DataType:}]
      \fort{MT}
\item [\textbf{Children:}]
      See \autoref{f:BCProperty}
\item [\textbf{Cardinality:}]
      0,1
\end{Ventryic}

\paragraph{\texttt{WallFunction\_t}}

A \fort{WallFunction\_t} node may be used as a child of
\fort{BCProperty\_t} to record data associated with the use of wall
function boundary conditions.

\textit{\uline{Node Attributes}}
\begin{Ventryic}{\textbf{Dimension Values:}}
\item [\textbf{Name:}]
      \fort{WallFunction}
\item [\textbf{Label:}]
      \fort{WallFunction\_t}
\item [\textbf{DataType:}]
      \fort{MT}
\item [\textbf{Children:}]
      See \autoref{f:WallFunction}
\item [\textbf{Cardinality:}]
      0,1
\end{Ventryic}

\paragraph{\texttt{Area\_t}}

An \fort{Area\_t} node may be used as a child of \fort{BCProperty\_t}
to record data associated with area-related boundary conditions such as
bleed.

\textit{\uline{Node Attributes}}
\begin{Ventryic}{\textbf{Dimension Values:}}
\item [\textbf{Name:}]
      \fort{Area}
\item [\textbf{Label:}]
      \fort{Area\_t}
\item [\textbf{DataType:}]
      \fort{MT}
\item [\textbf{Children:}]
      See \autoref{f:Area}
\item [\textbf{Cardinality:}]
      0,1
\end{Ventryic}

\subsubsection{Equation Specification Group}

Nodes in this group serve to identify the physical model associated
with the data being recorded. Nearly always, the data is of enumeration
type and is selected from a collection of terms defined in detail in
the SIDS.  The names are largely self explanatory, and the detailed
definitions will not be repeated here. Numerical values associated with
the physical model depend on the type of modeling being chosen and are
generally stored in child nodes of type \fort{DataArray\_t}.

\paragraph{\texttt{GoverningEquations\_t}}

This node names the equation set being solved, for example,
\fort{FullPotential} or \fort{NSTurbulent}. If Navier-Stokes, the
diffusion terms retained may be specified in a \fort{DiffusionModel}
subnode.

\textit{\uline{Node Attributes}}
\begin{Ventryic}{\textbf{Dimension Values:}}
\item [\textbf{Name:}]
      \fort{GoverningEquations}
\item [\textbf{Label:}]
      \fort{GoverningEquations\_t}
\item [\textbf{DataType:}]
      \fort{C1}
\item [\textbf{Dimension:}]
      1
\item [\textbf{Dimension Values:}]
      Length of string
\item [\textbf{Data:}]
      \fort{GoverningEquationsType} value
\item [\textbf{Children:}]
      See \autoref{f:GoverningEquations}
\item [\textbf{Cardinality:}]
      0,1
\item [\textbf{Parameters:}]
      \fort{CellDimension}
\end{Ventryic}

\paragraph{\texttt{GasModel\_t}}

A node of type \fort{GasModel\_t} names the gas model used, for example,
\fort{Ideal} or \fort{VanderWaals}.

\textit{\uline{Node Attributes}}
\begin{Ventryic}{\textbf{Dimension Values:}}
\item [\textbf{Name:}]
      \fort{GasModel}
\item [\textbf{Label:}]
      \fort{GasModel\_t}
\item [\textbf{DataType:}]
      \fort{C1}
\item [\textbf{Dimension:}]
      1
\item [\textbf{Dimension Values:}]
      Length of string
\item [\textbf{Data:}]
      \fort{GasModelType} value
\item [\textbf{Children:}]
      See \autoref{f:GasModel}
\item [\textbf{Cardinality:}]
      0,1
\end{Ventryic}

\paragraph{\texttt{ViscosityModel\_t}}

A node of type \fort{ViscosityModel\_t} names the molecular viscosity
model used to relate the viscosity to the temperature, for example,
\fort{PowerLaw} or \fort{SutherlandLaw}.

\textit{\uline{Node Attributes}}
\begin{Ventryic}{\textbf{Dimension Values:}}
\item [\textbf{Name:}]
      \fort{ViscosityModel}
\item [\textbf{Label:}]
      \fort{ViscosityModel\_t}
\item [\textbf{DataType:}]
      \fort{C1}
\item [\textbf{Dimension:}]
      1
\item [\textbf{Dimension Values:}]
      Length of string
\item [\textbf{Data:}]
      \fort{ViscosityModelType} value
\item [\textbf{Children:}]
      See \autoref{f:ViscosityModel}
\item [\textbf{Cardinality:}]
      0,1
\end{Ventryic}

\paragraph{\texttt{EquationDimension}}

A node named \fort{EquationDimension}, of type \fort{int[]}, gives the
number of dependent variables required for a complete solution
description, or the number of equations being solved. For example, for
\fort{NSTurbulent} with the $k$-$\epsilon$ turbulence model in three
dimensions, it is 7.

\textit{\uline{Node Attributes}}
\begin{Ventryic}{\textbf{Dimension Values:}}
\item [\textbf{Name:}]
      \fort{EquationDimension}
\item [\textbf{Label:}]
      ``\fort{int}''
\item [\textbf{DataType:}]
      \fort{I4}
\item [\textbf{Dimension:}]
      1
\item [\textbf{Dimension Values:}]
      1
\item [\textbf{Data:}]
      \fort{EquationDimension} value
\item [\textbf{Children:}]
      None
\item [\textbf{Cardinality:}]
      0,1
\end{Ventryic}

\paragraph{\texttt{ThermalConductivityModel\_t}}

A node of type \fort{ThermalConductivityModel\_t} names the model used
to relate the thermal conductivity to the temperature, for example,
\fort{ConstantPrandtl}, \fort{PowerLaw}, or \fort{SutherlandLaw}. These
closely parallel the viscosity model.

\textit{\uline{Node Attributes}}
\begin{Ventryic}{\textbf{Dimension Values:}}
\item [\textbf{Name:}]
      \fort{ThermalConductivityModel}
\item [\textbf{Label:}]
      \fort{ThermalConductivityModel\_t}
\item [\textbf{DataType:}]
      \fort{C1}
\item [\textbf{Dimension:}]
      1
\item [\textbf{Dimension Values:}]
      Length of string
\item [\textbf{Data:}]
      \fort{ThermalConductivityModelType} value
\item [\textbf{Children:}]
      See \autoref{f:ThermalConductivityModel}
\item [\textbf{Cardinality:}]
      0,1
\end{Ventryic}

\paragraph{\texttt{TurbulenceClosure\_t}}

\enlargethispage{\baselineskip}%
A node of type \fort{TurbulenceClosure\_t} names the method of closing
the Reynolds stress equations when the governing equations are
turbulent, for example, \fort{EddyViscosity} or
\fort{ReynoldsStressAlgebraic}.

\textit{\uline{Node Attributes}}
\begin{Ventryic}{\textbf{Dimension Values:}}
\item [\textbf{Name:}]
      \fort{TurbulenceClosure}
\item [\textbf{Label:}]
      \fort{TurbulenceClosure\_t}
\item [\textbf{DataType:}]
      \fort{C1}
\item [\textbf{Dimension:}]
      1
\item [\textbf{Dimension Values:}]
      Length of string
\item [\textbf{Data:}]
      \fort{TurbulenceClosureType} value
\item [\textbf{Children:}]
      See \autoref{f:TurbulenceClosure}
\item [\textbf{Cardinality:}]
      0,1
\end{Ventryic}

\paragraph{\texttt{TurbulenceModel\_t}}

A node of type \fort{TurbulenceModel\_t} names the equation
set used to model the turbulence quantities, for example,
\fort{Algebraic\_Bald\-winLomax} or \fort{OneEquation\_Spa\-lartAll\-maras}.

\textit{\uline{Node Attributes}}
\begin{Ventryic}{\textbf{Dimension Values:}}
\item [\textbf{Name:}]
      \fort{TurbulenceModel}
\item [\textbf{Label:}]
      \fort{TurbulenceModel\_t}
\item [\textbf{DataType:}]
      \fort{C1}
\item [\textbf{Dimension:}]
      1
\item [\textbf{Dimension Values:}]
      Length of string
\item [\textbf{Data:}]
      \fort{TurbulenceModelType} value
\item [\textbf{Children:}]
      See \autoref{f:TurbulenceModel}
\item [\textbf{Cardinality:}]
      0,1
\item [\textbf{Parameters:}]
      \fort{CellDimension}
\end{Ventryic}

\paragraph{\texttt{ThermalRelaxationModel\_t}}

A node of type \fort{ThermalRelaxationModel\_t} names the equation
set used to model the thermal relaxation quantities, for example,
\fort{Frozen} or \fort{ThermalEquilib}.

\textit{\uline{Node Attributes}}
\begin{Ventryic}{\textbf{Dimension Values:}}
\item [\textbf{Name:}]
      \fort{ThermalRelaxationModel}
\item [\textbf{Label:}]
      \fort{ThermalRelaxationModel\_t}
\item [\textbf{DataType:}]
      \fort{C1}
\item [\textbf{Dimension:}]
      1
\item [\textbf{Dimension Values:}]
      Length of string
\item [\textbf{Data:}]
      \fort{ThermalRelaxationModelType} value
\item [\textbf{Children:}]
      See \autoref{f:ThermalRelaxationModel}
\item [\textbf{Cardinality:}]
      0,1
\end{Ventryic}

\paragraph{\texttt{ChemicalKineticsModel\_t}}

A node of type \fort{ChemicalKineticsModel\_t} names the equation
set used to model the chemical kinetics quantities, for example,
\fort{Frozen} or \fort{ChemicalEquilibCurveFit}.

\textit{\uline{Node Attributes}}
\begin{Ventryic}{\textbf{Dimension Values:}}
\item [\textbf{Name:}]
      \fort{ChemicalKineticsModel}
\item [\textbf{Label:}]
      \fort{ChemicalKineticsModel\_t}
\item [\textbf{DataType:}]
      \fort{C1}
\item [\textbf{Dimension:}]
      1
\item [\textbf{Dimension Values:}]
      Length of string
\item [\textbf{Data:}]
      \fort{ChemicalKineticsModelType} value
\item [\textbf{Children:}]
      See \autoref{f:ChemicalKineticsModel}
\item [\textbf{Cardinality:}]
      0,1
\end{Ventryic}

\paragraph{\texttt{EMElectricFieldModel\_t}}

A node of type \fort{EMElectricFieldModel\_t} names the electric
field model used for electromagnetic flows, for example,
\fort{Constant} or \fort{Voltage}.

\textit{\uline{Node Attributes}}
\begin{Ventryic}{\textbf{Dimension Values:}}
\item [\textbf{Name:}]
      \fort{EMElectricFieldModel}
\item [\textbf{Label:}]
      \fort{EMElectricFieldModel\_t}
\item [\textbf{DataType:}]
      \fort{C1}
\item [\textbf{Dimension:}]
      1
\item [\textbf{Dimension Values:}]
      Length of string
\item [\textbf{Data:}]
      \fort{EMElectricFieldModelType} value
\item [\textbf{Children:}]
      See \autoref{f:EMElectricFieldModel}
\item [\textbf{Cardinality:}]
      0,1
\end{Ventryic}

\paragraph{\texttt{EMMagneticFieldModel\_t}}

A node of type \fort{EMMagneticFieldModel\_t} names the magnetic
field model used for electromagnetic flows, for example,
\fort{Constant} or \fort{Interpolated}.

\textit{\uline{Node Attributes}}
\begin{Ventryic}{\textbf{Dimension Values:}}
\item [\textbf{Name:}]
      \fort{EMMagneticFieldModel}
\item [\textbf{Label:}]
      \fort{EMMagneticFieldModel\_t}
\item [\textbf{DataType:}]
      \fort{C1}
\item [\textbf{Dimension:}]
      1
\item [\textbf{Dimension Values:}]
      Length of string
\item [\textbf{Data:}]
      \fort{EMMagneticFieldModelType} value
\item [\textbf{Children:}]
      See \autoref{f:EMMagneticFieldModel}
\item [\textbf{Cardinality:}]
      0,1
\end{Ventryic}

\paragraph{\texttt{EMConductivityModel\_t}}

A node of type \fort{EMConductivityModel\_t} names the conductivity
model used for electromagnetic flows, for example,
\fort{Constant} or \fort{Equilibrium\_LinRessler}.

\textit{\uline{Node Attributes}}
\begin{Ventryic}{\textbf{Dimension Values:}}
\item [\textbf{Name:}]
      \fort{EMConductivityModel}
\item [\textbf{Label:}]
      \fort{EMConductivityModel\_t}
\item [\textbf{DataType:}]
      \fort{C1}
\item [\textbf{Dimension:}]
      1
\item [\textbf{Dimension Values:}]
      Length of string
\item [\textbf{Data:}]
      \fort{EMConductivityModelType} value
\item [\textbf{Children:}]
      See \autoref{f:EMConductivityModel}
\item [\textbf{Cardinality:}]
      0,1
\end{Ventryic}

\paragraph{\texttt{FlowEquationSet\_t}}

A node of type \fort{FlowEquationSet\_t} appears either at the highest
level of the tree (under \fort{CGNSBase\_t}), to indicate the equation
set whose solution is recorded throughout the database, or below a
\fort{Zone\_t} node, to indicate the set of equations solved in that
zone. The usual convention applies, i.e., specifications at the local
(zone) level override global specifications.

\textit{\uline{Node Attributes}}
\begin{Ventryic}{\textbf{Dimension Values:}}
\item [\textbf{Name:}]
      \fort{FlowEquationSet}
\item [\textbf{Label:}]
      \fort{FlowEquationSet\_t}
\item [\textbf{DataType:}]
      \fort{MT}
\item [\textbf{Children:}]
      See \autoref{f:FlowEquationSet}
\item [\textbf{Cardinality:}]
      0,1
\item [\textbf{Parameters:}]
      \fort{CellDimension}
\end{Ventryic}

\subsubsection{Family Group}

Because there is rarely a 1-to-1 connection between mesh regions and
geometric entities, it is often desirable to set geometric associations
indirectly in a CGNS file.
That is, rather than setting the geometry data for each mesh entity
(nodes, edges, and faces), it's useful to associate them with
intermediate objects.
The intermediate objects are in turn linked to nodal regions of the
computational mesh.
This intermediate object is defined as a \emph{CFD family}.

Each mesh surface may linked to the geometric entities of one or more CAD
databases by a user-defined CFD family name.
The CFD family corresponds to one or more CAD geometric entities on which
the mesh face is projected.
Each one of these geometric entities is described in a CAD file and is
not redefined within the CGNS file.

\paragraph{\texttt{Family\_t}}

This node, a child of the \fort{CGNSBase\_t} node, contains the
definition of a single CFD family.
Multiple \texttt{Family\_t} nodes are allowed.

\textit{\uline{Node Attributes}}
\begin{Ventryic}{\textbf{Dimension Values:}}
\item [\textbf{Name:}]
      User defined
\item [\textbf{Label:}]
      \fort{Family\_t}
\item [\textbf{DataType:}]
      \fort{MT}
\item [\textbf{Children:}]
      See \autoref{f:Family}
\item [\textbf{Cardinality:}]
      0,$N$
\end{Ventryic}

\paragraph{\texttt{FamilyName\_t}}

This node is used to identify a family to which
a particular zone or boundary belongs.
Note that the name of the family is defined by the ``Name''
of the \fort{Family\_t} node, and is stored as data in the
\fort{FamilyName\_t} node.

\textit{\uline{Node Attributes}}
\begin{Ventryic}{\textbf{Dimension Values:}}
\item [\textbf{Name:}]
      \fort{FamilyName}
\item [\textbf{Label:}]
      \fort{FamilyName\_t}
\item [\textbf{DataType:}]
      \fort{C1}
\item [\textbf{Dimension:}]
      1
\item [\textbf{Dimension Values:}]
      Length of string
\item [\textbf{Data:}]
      Name of CFD family
\item [\textbf{Children:}]
      None
\item [\textbf{Cardinality:}]
      0,1
\end{Ventryic}

\paragraph{\texttt{FamilyBC\_t}}

This node contains a boundary condition type for a
particular CFD family.

\textit{\uline{Node Attributes}}
\begin{Ventryic}{\textbf{Dimension Values:}}
\item [\textbf{Name:}]
      \fort{FamilyBC}
\item [\textbf{Label:}]
      \fort{FamilyBC\_t}
\item [\textbf{DataType:}]
      \fort{C1}
\item [\textbf{Dimension:}]
      1
\item [\textbf{Dimension Values:}]
      Length of string
\item [\textbf{Data:}]
      \fort{BCType} value
\item [\textbf{Children:}]
      See \autoref{f:FamilyBC}
\item [\textbf{Cardinality:}]
      0,1
\end{Ventryic}

\paragraph{\texttt{FamilyBCDataSet\_t}}

A \fort{FamilyBCDataSet\_t} node contains Dirichlet and Neumann
data for a single set of boundary condition equations to be applied
to a Family. It's intended use is for simple boundary condition types,
where the equations imposed do not depend on local flow conditions.

\textit{\uline{Node Attributes}}
\begin{Ventryic}{\textbf{Dimension Values:}}
\item [\textbf{Name:}]
      User defined
\item [\textbf{Label:}]
      \fort{FamilyBCDataSet\_t}
\item [\textbf{DataType:}]
      \fort{C1}
\item [\textbf{Dimension:}]
      1
\item [\textbf{Dimension Values:}]
      Length of string
\item [\textbf{Data:}]
      \fort{BCTypeSimple} value
\item [\textbf{Children:}]
      See \autoref{f:FamilyBCDataSet}
\item [\textbf{Cardinality:}]
      0,$N$
\end{Ventryic}

\paragraph{\texttt{GeometryReference\_t}}

\texttt{GeometryReference\_t} nodes are used to associate
a CFD family with one or more CAD databases.

\textit{\uline{Node Attributes}}
\begin{Ventryic}{\textbf{Dimension Values:}}
\item [\textbf{Name:}]
      User defined
\item [\textbf{Label:}]
      \fort{GeometryReference\_t}
\item [\textbf{DataType:}]
      \fort{MT}
\item [\textbf{Children:}]
      See \autoref{f:GeometryReference}
\item [\textbf{Cardinality:}]
      0,$N$
\end{Ventryic}

\paragraph{\texttt{GeometryFile\_t}}

This node contains the name of the CAD geometry file.

\textit{\uline{Node Attributes}}
\begin{Ventryic}{\textbf{Dimension Values:}}
\item [\textbf{Name:}]
      \fort{GeometryFile}
\item [\textbf{Label:}]
      \fort{GeometryFile\_t}
\item [\textbf{DataType:}]
      \fort{C1}
\item [\textbf{Dimension:}]
      1
\item [\textbf{Dimension Values:}]
      Length of string
\item [\textbf{Data:}]
      Name of the CAD geometry file
\item [\textbf{Children:}]
      None
\item [\textbf{Cardinality:}]
      1
\end{Ventryic}

\paragraph{\texttt{GeometryFormat\_t}}

This enumeration node defines the format of the CAD geometry file.

\textit{\uline{Node Attributes}}
\begin{Ventryic}{\textbf{Dimension Values:}}
\item [\textbf{Name:}]
      \fort{GeometryFormat}
\item [\textbf{Label:}]
      \fort{GeometryFormat\_t}
\item [\textbf{DataType:}]
      \fort{C1}
\item [\textbf{Dimension:}]
      1
\item [\textbf{Dimension Values:}]
      Length of string
\item [\textbf{Data:}]
      Name of the CAD geometry format
\item [\textbf{Children:}]
      None
\item [\textbf{Cardinality:}]
      1
\end{Ventryic}

\paragraph{\texttt{GeometryEntity\_t}}

\texttt{GeometryEntity\_t} nodes define the names of the
entities in CAD geometry file that make up a CFD family.

\textit{\uline{Node Attributes}}
\begin{Ventryic}{\textbf{Dimension Values:}}
\item [\textbf{Name:}]
      User defined
\item [\textbf{Label:}]
      \fort{GeometryEntity\_t}
\item [\textbf{DataType:}]
      \fort{MT}
\item [\textbf{Children:}]
      None
\item [\textbf{Cardinality:}]
      0,$N$
\end{Ventryic}

\subsubsection{Time-Dependent Group}
\label{s:timedep}

Nodes in this section are used for information related to time-dependent
flows, and include specification of grid motion and storage of
time-dependent or iterative data.

\paragraph{\texttt{BaseIterativeData\_t}}

Located directly under the \fort{CGNSBase\_t} node, the
\fort{BaseIterativeData\_t} node contains information about the number
of time steps or iterations being recorded, and the time and/or
iteration values at each step.
In addition, it may include the list of zones and families for each step
of the simulation, if these vary throughout the simulation.

\textit{\uline{Node Attributes}}
\begin{Ventryic}{\textbf{Dimension Values:}}
\item [\textbf{Name:}]
      User defined
\item [\textbf{Label:}]
      \fort{BaseIterativeData\_t}
\item [\textbf{DataType:}]
      \fort{I4}
\item [\textbf{Dimension:}]
      1
\item [\textbf{Dimension Values:}]
      1
\item [\textbf{Data:}]
      \fort{NumberOfSteps}
\item [\textbf{Children:}]
      See \autoref{f:BaseIterativeData}
\item [\textbf{Cardinality:}]
      0,1
\end{Ventryic}

\paragraph{\texttt{ZoneIterativeData\_t}}

The \fort{ZoneIterativeData\_t} node is a child of the \fort{Zone\_t}
node, and is used to store pointers to zonal data for each recorded
step of the simulation.

\textit{\uline{Node Attributes}}
\begin{Ventryic}{\textbf{Dimension Values:}}
\item [\textbf{Name:}]
      User defined
\item [\textbf{Label:}]
      \fort{ZoneIterativeData\_t}
\item [\textbf{DataType:}]
      \fort{MT}
\item [\textbf{Children:}]
      See \autoref{f:ZoneIterativeData}
\item [\textbf{Cardinality:}]
      0,1
\item [\textbf{Parameters:}]
      \fort{NumberOfSteps}
\end{Ventryic}

\paragraph{\texttt{RigidGridMotion\_t}}

\fort{RigidGridMotion\_t} nodes are used to store data
defining rigid translation and/or rotation of the grid coordinates.
Multiple \fort{RigidGridMotion\_t} nodes may be associated with different
iterations or time steps in the computation.
This association is recorded under the \fort{ZoneIterativeData\_t} node.

\textit{\uline{Node Attributes}}
\begin{Ventryic}{\textbf{Dimension Values:}}
\item [\textbf{Name:}]
      User defined
\item [\textbf{Label:}]
      \fort{RigidGridMotion\_t}
\item [\textbf{DataType:}]
      \fort{C1}
\item [\textbf{Dimension:}]
      1
\item [\textbf{Dimension Values:}]
      Length of string
\item [\textbf{Data:}]
      \fort{RigidGridMotionType} value
\item [\textbf{Children:}]
      See \autoref{f:RigidGridMotion}
\item [\textbf{Cardinality:}]
      0,$N$
\end{Ventryic}

\paragraph{\texttt{ArbitraryGridMotion\_t}}

\fort{ArbitraryGridMotion\_t} nodes are used to store grid
velocities for each grid point in a zone (i.e., for deforming grids).
Multiple \fort{ArbitraryGridMotion\_t} nodes may be associated with
different iterations or time steps in the computation.
This association is recorded under the \fort{ZoneIterativeData\_t} node.

Note that instantaneous grid coordinates at different iterations or time
steps may be recorded using multiple \fort{GridCoordinates\_t} nodes
under a \fort{Zone\_t} node.

\textit{\uline{Node Attributes}}
\begin{Ventryic}{\textbf{Dimension Values:}}
\item [\textbf{Name:}]
      User defined
\item [\textbf{Label:}]
      \fort{ArbitraryGridMotion\_t}
\item [\textbf{DataType:}]
      \fort{C1}
\item [\textbf{Dimension:}]
      1
\item [\textbf{Dimension Values:}]
      Length of string
\item [\textbf{Data:}]
      \fort{ArbitraryGridMotionType} value
\item [\textbf{Children:}]
      See \autoref{f:ArbitraryGridMotion}
\item [\textbf{Cardinality:}]
      0,$N$
\item [\textbf{Parameters:}]
      \fort{IndexDimension}, \fort{VertexSize}, \fort{CellSize}
\item [\textbf{Functions:}]
      \fort{DataSize}
\end{Ventryic}

\subsubsection{Structural Nodes}

In this section we describe the highest levels of the hierarchy. Nodes
in this section store only quantities which refer to all the entities
below them. Therir primary function is to provide organization to the
data below.

\paragraph{\texttt{Zone\_t}}
\label{s:Zone}

Directly below the highest level node in the database, which is
by definition of type \fort{CGNSBase\_t}, are found nodes of type
\fort{Zone\_t} providing entry into the data specific to each
zone. There are as many \fort{Zone\_t} nodes as there are zones. Their
children, in turn, record grid, field, connectivity, and boundary
conditions, and a variety of auxiliary data.

\textit{\uline{Node Attributes}}
\begin{Ventryic}{\textbf{Dimension Values:}}
\item [\textbf{Name:}]
      User defined
\item [\textbf{Label:}]
      \fort{Zone\_t}
\item [\textbf{DataType:}]
      \ital{cgsize\_t}
\item [\textbf{Dimension:}]
      2
\item [\textbf{Dimension Values:}]
      \fort{IndexDimension}, 3
\item [\textbf{Data:}]
      \fort{VertexSize[IndexDimension]}, \fort{CellSize[IndexDimension]},
      \fort{VertexSizeBoundary[IndexDimension]}
\item [\textbf{Children:}]
      See \autoref{f:Zone}
\item [\textbf{Cardinality:}]
      0,$N$
\item [\textbf{Parameters:}]
      \fort{CellDimension}, \fort{PhysicalDimension}
\end{Ventryic}

\paragraph{\texttt{CGNSBase\_t}}

The \fort{CGNSBase\_t} node is by definition the highest level node in
the database, and is located directly below the ADF root node.
It provides entry into all other data.
Multiple \fort{CGNSBase\_t} nodes are allowed in an ADF file.
The particular database being accessed is determined by the name
of the \fort{CGNSBase\_t} node.

The only data stored in the node itself are \fort{CellDimension}, the
dimensionality of a cell in the mesh (i.e., 3 for a volume cell and 2
for a face cell), and \fort{PhysicalDimension}, the number of indices
required to specify a unique physical location in the field data being
recorded. However, a variety of global information concerning the entire
database may be stored in children of the \fort{CGNSBase\_t} node. In
particular, a \fort{Descriptor\_t} node at this level can store user
commentary on the entire history of the development of the database.

Other information typically stored directly below the \fort{CGNSBase\_t}
node includes convergence histories, reference states, dimensional
units, integrated quantities, and information on the flow equations
being solved.

\textit{\uline{Node Attributes}}
\begin{Ventryic}{\textbf{Dimension Values:}}
\item [\textbf{Name:}]
      User defined
\item [\textbf{Label:}]
      \fort{CGNSBase\_t}
\item [\textbf{DataType:}]
      \fort{I4}
\item [\textbf{Dimension:}]
      1
\item [\textbf{Dimension Values:}]
      2
\item [\textbf{Data:}]
      \fort{CellDimension}, \fort{PhysicalDimension}
\item [\textbf{Children:}]
      See \autoref{f:CGNSBase}
\item [\textbf{Cardinality:}]
      0,$N$
\end{Ventryic}

\paragraph{\texttt{SimulationType\_t}}

This enumeration-type node is a child of the \fort{CGNSBase\_t} node,
and specifies whether or not the data below \fort{CGNSBase\_t} is
time-accurate.
Nodes for describing time-dependent data are presented in
\autoref{s:timedep}.

\textit{\uline{Node Attributes}}
\begin{Ventryic}{\textbf{Dimension Values:}}
\item [\textbf{Name:}]
      \fort{SimulationType}
\item [\textbf{Label:}]
      \fort{SimulationType\_t}
\item [\textbf{DataType:}]
      \fort{C1}
\item [\textbf{Dimension:}]
      1
\item [\textbf{Dimension Values:}]
      Length of string
\item [\textbf{Data:}]
      \fort{TimeAccurate} or \fort{NonTimeAccurate}
\item [\textbf{Children:}]
      None
\item [\textbf{Cardinality:}]
      0,1
\end{Ventryic}

\paragraph{\texttt{ZoneType\_t}}

This enumeration-type node is a required child of the \fort{Zone\_t}
node, and specifies whether the grid in that zone is structured or
unstructured.

\textit{\uline{Node Attributes}}
\begin{Ventryic}{\textbf{Dimension Values:}}
\item [\textbf{Name:}]
      \fort{ZoneType}
\item [\textbf{Label:}]
      \fort{ZoneType\_t}
\item [\textbf{DataType:}]
      \fort{C1}
\item [\textbf{Dimension:}]
      1
\item [\textbf{Dimension Values:}]
      Length of string
\item [\textbf{Data:}]
      \fort{Structured} or \fort{Unstructured}
\item [\textbf{Children:}]
      None
\item [\textbf{Cardinality:}]
      1
\end{Ventryic}

\paragraph{\texttt{CGNSLibraryVersion\_t}}

An ADF file containing a CGNS database also contains, directly below the
ADF root node, a \fort{CGNSLibraryVersion\_t} node.
This node contains the version number of the CGNS standard with which
the file is consistent, and is created automatically when the file is
created or modified using the CGNS Mid-Level Library.
Note that this node is not actually part of the CGNS database, since
it is not located below a \fort{CGNSBase\_t} node.

Note that an ADF file may contain multiple CGNS databases, but there is
only one \fort{CGNSLibraryVersion\_t} node.
It is assumed that the version number in the \fort{CGNSLibraryVersion\_t}
node is applicable to all the CGNS databases in the file.

Note also that some CGNS nodes may actually be links to CGNS nodes in
other files.
In this case, it is assumed that the \fort{CGNSLibraryVersion\_t} node
in the ``top-level'' file is applicable to the file(s) containing the
linked nodes.

\textit{\uline{Node Attributes}}
\begin{Ventryic}{\textbf{Dimension Values:}}
\item [\textbf{Name:}]
      \fort{CGNSLibraryVersion}
\item [\textbf{Label:}]
      \fort{CGNSLibraryVersion\_t}
\item [\textbf{DataType:}]
      \fort{R4}
\item [\textbf{Dimension:}]
      1
\item [\textbf{Dimension Values:}]
      1
\item [\textbf{Data:}]
      CGNS version number
\item [\textbf{Children:}]
      None
\item [\textbf{Cardinality:}]
      1
\end{Ventryic}


\clearemptydoublepage
\appendix
%\fancyhead[RO]{\bfseries \rightmark}
%\app{CGNS File Mapping Figures}
\section{CGNS File Mapping Figures}
\label{s:figures}
\thispagestyle{plain}

This Appendix contains figures illustrating the nodal structure of
a CGNS database.
An example HDF5 file hierarchy is shown in \autoref{f:example_hierarchy},
but without any internal detail of the individual nodes.
Just below the HDF5 root node is a \fort{CGNSLibraryVersion\_t} node, and
one or more \fort{CGNSBase\_t} nodes.
Each \fort{CGNSBase\_t} node represents the root node of a CGNS database.

\autoref{f:example_node} shows the layout of succeeding figures.
At the top of each page is a single parent node, and below it are boxes
containing detailed descriptions of the nodes which the mapping specifies
as children.
Also within the parent's box are references to the figures corresponding
to its parents in turn.

\renewcommand{\topfraction}{1.00}
\renewcommand{\bottomfraction}{1.00}
\renewcommand{\textfraction}{0.0}
\setlength{\unitlength}{1em}	% unit length for example hierarchy figure

\newlength{\indName}
\newlength{\indLabel}
\newlength{\indDimensions}
\newlength{\indDimensionValues}
\newlength{\indData}
\newlength{\indParameters}
\newlength{\indFunctions}
\newlength{\indChildNodes}
\newlength{\indbullet}
\newlength{\spacea}
\newlength{\spaceb}
\newlength{\spacec}
\settowidth{\indName}{\textbf{Name:} }
\settowidth{\indLabel}{\textbf{Label:} }
\settowidth{\indDimensions}{\textbf{Dimensions:} }
\settowidth{\indDimensionValues}{\textbf{Dimension Values:} }
\settowidth{\indData}{\textbf{Data:} }
\settowidth{\indParameters}{\textbf{Parameters:} }
\settowidth{\indFunctions}{\textbf{Functions:} }
\settowidth{\indChildNodes}{\textbf{Child Nodes:} }
\settowidth{\indbullet}{\textbullet ~~}
\setlength{\spacea}{\indDimensionValues-\indChildNodes}
\setlength{\spaceb}{\indDimensions-\indParameters+5em}
\setlength{\spacec}{\indDimensions-\indFunctions+5em}

\newlength{\textmcap}
\setlength{\textmcap}{\textheight-\abovecaptionskip-\baselineskip}

\begin{landscape}
\begin{figure}[!htp]
\centering
% In example_hierarchy.pic, total width of picture changed from 71 to
% the correct value of 64.  (Is there a bug in flow?)
\resizebox{\linewidth}{!}{\input{figures/example_hierarchy_cet.pic}}
\caption{Example Hierarchy}
\label{f:example_hierarchy}
\end{figure}
\end{landscape}

\newpage
\setlength{\unitlength}{\baselineskip}	% unit length for remaining figures
\begin{figure}[!htp]
\centering
\resizebox{\linewidth}{!}{\input{figures/example_node_cet.pic}}
\caption{Example Node Structure}
\label{f:example_node}
\end{figure}

\clearpage
\begin{figure}[!htp]
\centering
\resizebox{\linewidth}{!}{\input{figures/CGNSBase1.pic}}
\caption[\fort{CGNSBase\_t} Data Structure]{\fort{CGNSBase\_t} Data Structure (\textit{Continued on next page})}
\label{f:CGNSBase}
\end{figure}
\clearpage
\begin{figure}[!htp]
\centering
\resizebox{\linewidth}{!}{\input{figures/CGNSBase2.pic}}
\par\vspace{\abovecaptionskip}
\textbf{Figure \arabic{figure}:} \fort{CGNSBase\_t} Data Structure (\textit{Continued from previous page})
\end{figure}

\clearpage
\begin{figure}[!htp]
\centering
\resizebox{\linewidth}{!}{\input{figures/Zone1.pic}}
\caption[\fort{Zone\_t} Data Structure]{\fort{Zone\_t} Data Structure (\textit{Continued on next page})}
\label{f:Zone}
\end{figure}
\clearpage
\begin{figure}[!htp]
\centering
\resizebox{\linewidth}{!}{\input{figures/Zone2.pic}}
\par\vspace{\abovecaptionskip}
\textbf{Figure \arabic{figure}:} \fort{Zone\_t} Data Structure (\textit{Continued from previous page})
\end{figure}

\clearpage
\begin{figure}[!htp]
\centering
\input{figures/GridCoordinates.pic}
\caption{\fort{GridCoordinates\_t} Data Structure}
\label{f:GridCoordinates}
\end{figure}

\clearpage
\begin{figure}[!htp]
\centering
\input{figures/Elements.pic}
\caption{\fort{Elements\_t} Data Structure}
\label{f:Elements}
\end{figure}

\clearpage
\begin{figure}[!htp]
\centering
%\input{figures/Axisymmetry.pic}
\resizebox{\linewidth}{!}{\input{figures/Axisymmetry.pic}}
\caption{\fort{Axisymmetry\_t} Data Structure}
\label{f:Axisymmetry}
\end{figure}

\clearpage
\begin{figure}[!htp]
\centering
\input{figures/RotatingCoordinates.pic}
\caption{\fort{RotatingCoordinates\_t} Data Structure}
\label{f:RotatingCoordinates}
\end{figure}

\clearpage
\begin{figure}[!htp]
\centering
\resizebox{!}{\textheight-2\baselineskip}{\input{figures/FlowSolution.pic}}
%\input{figures/FlowSolution.pic}
\caption{\fort{FlowSolution\_t} and \fort{DiscreteData\_t} Data Structures}
\label{f:FlowSolution}
\end{figure}

\clearpage
\begin{figure}[!htp]
\centering
\resizebox{!}{\textheight-2\baselineskip}{\input{figures/ZoneGridConnectivity.pic}}
%\input{figures/ZoneGridConnectivity.pic}
\caption{\fort{ZoneGridConnectivity\_t} Data Structure}
\label{f:ZoneGridConnectivity}
\end{figure}

\clearpage
\begin{figure}[!htp]
\centering
\resizebox{!}{\textheight-2\baselineskip}{\input{figures/GridConnectivity1to1.pic}}
%\input{figures/GridConnectivity1to1.pic}
\caption{\fort{GridConnectivity1to1\_t} Data Structure}
\label{f:GridConnectivity1to1}
\end{figure}

\clearpage
\begin{figure}[!htp]
\centering
\resizebox{\linewidth}{!}{\input{figures/GridConnectivity.pic}}
\caption{\fort{GridConnectivity\_t} Data Structure}
\label{f:GridConnectivity}
\end{figure}

\clearpage
\begin{figure}[!htp]
\centering
\input{figures/GridConnectivityProperty.pic}
\caption{\fort{GridConnectivityProperty\_t} Data Structure}
\label{f:GridConnectivityProperty}
\end{figure}

\clearpage
\begin{figure}[!htp]
\centering
%\input{figures/Periodic.pic}
\resizebox{\linewidth}{!}{\input{figures/Periodic.pic}}
\caption{\fort{Periodic\_t} Data Structure}
\label{f:Periodic}
\end{figure}

\clearpage
\begin{figure}[!htp]
\centering
\input{figures/AverageInterface.pic}
\caption{\fort{AverageInterface\_t} Data Structure}
\label{f:AverageInterface}
\end{figure}

\clearpage
\begin{figure}[!htp]
\centering
\input{figures/OversetHoles.pic}
\caption{\fort{OversetHoles\_t} Data Structure}
\label{f:OversetHoles}
\end{figure}

\clearpage
\begin{figure}[!htp]
\centering
\input{figures/ZoneBC.pic}
\caption{\fort{ZoneBC\_t} Data Structure}
\label{f:ZoneBC}
\end{figure}

%\clearpage
%\begin{figure}[!htp]
%\centering
%%\resizebox{\linewidth}{!}{\input{figures/BC.pic}}
%\resizebox{!}{\textheight-2\baselineskip}{\input{figures/BC.pic}}
%\caption{\fort{BC\_t} Data Structure}
%\label{f:BC}
%\end{figure}

\clearpage
\begin{figure}[!htp]
\centering
\resizebox{\linewidth}{!}{\input{figures/BC1.pic}}
\caption[\fort{BC\_t} Data Structure]{\fort{BC\_t} Data Structure (\textit{Continued on next page})}
\label{f:BC}
\end{figure}
\clearpage
\begin{figure}[!htp]
\centering
\resizebox{\linewidth}{!}{\input{figures/BC2.pic}}
\par\vspace{\abovecaptionskip}
\textbf{Figure \arabic{figure}:} \fort{BC\_t} Data Structure (\textit{Continued from previous page})
\end{figure}

\clearpage
\begin{figure}[!htp]
\centering
%\input{figures/BCDataSet.pic}
\resizebox{!}{\textheight-2\baselineskip}{\input{figures/BCDataSet.pic}}
\caption{\fort{BCDataSet\_t} Data Structure}
\label{f:BCDataSet}
\end{figure}

\clearpage
\begin{figure}[!htp]
\centering
\input{figures/BCData.pic}
\caption{\fort{BCData\_t} Data Structure}
\label{f:BCData}
\end{figure}

\clearpage
\begin{figure}[!htp]
\centering
\input{figures/BCProperty.pic}
\caption{\fort{BCProperty\_t} Data Structure}
\label{f:BCProperty}
\end{figure}

\clearpage
\begin{figure}[!htp]
\centering
\input{figures/WallFunction.pic}
\caption{\fort{WallFunction\_t} Data Structure}
\label{f:WallFunction}
\end{figure}

\clearpage
\begin{figure}[!htp]
\centering
\input{figures/Area.pic}
\caption{\fort{Area\_t} Data Structure}
\label{f:Area}
\end{figure}

%\clearpage
%\begin{figure}[!htp]
%\centering
%%\input{figures/FlowEquationSet.pic}
%%\resizebox{\linewidth}{!}{\input{figures/FlowEquationSet.pic}}
%\resizebox{!}{\textheight-2\baselineskip}{\input{figures/FlowEquationSet.pic}}
%\caption{\fort{FlowEquationSet\_t} Data Structure}
%\label{f:FlowEquationSet}
%\end{figure}

\clearpage
\begin{figure}[!htp]
\centering
\resizebox{\linewidth}{!}{\input{figures/FlowEquationSet1.pic}}
\caption[\fort{FlowEquationSet\_t} Data Structure]{\fort{FlowEquationSet\_t} Data Structure (\textit{Continued on next page})}
\label{f:FlowEquationSet}
\end{figure}
\clearpage
\begin{figure}[!htp]
\centering
\resizebox{\linewidth}{!}{\input{figures/FlowEquationSet2.pic}}
\par\vspace{\abovecaptionskip}
\textbf{Figure \arabic{figure}:} \fort{FlowEquationSet\_t} Data Structure (\textit{Continued from previous page})
\end{figure}

\clearpage
\begin{figure}[!htp]
\centering
\input{figures/GoverningEquations.pic}
\caption{\fort{GoverningEquations\_t} Data Structure}
\label{f:GoverningEquations}
\end{figure}

\clearpage
\begin{figure}[!htp]
\centering
\input{figures/GasModel.pic}
\caption{\fort{GasModel\_t} Data Structure}
\label{f:GasModel}
\end{figure}

\clearpage
\begin{figure}[!htp]
\centering
\input{figures/ViscosityModel.pic}
\caption{\fort{ViscosityModel\_t} Data Structure}
\label{f:ViscosityModel}
\end{figure}

\clearpage
\begin{figure}[!htp]
\centering
\input{figures/ThermalConductivityModel.pic}
\caption{\fort{ThermalConductivityModel\_t} Data Structure}
\label{f:ThermalConductivityModel}
\end{figure}

\clearpage
\begin{figure}[!htp]
\centering
\input{figures/TurbulenceClosure.pic}
\caption{\fort{TurbulenceClosure\_t} Data Structure}
\label{f:TurbulenceClosure}
\end{figure}

\clearpage
\begin{figure}[!htp]
\centering
\input{figures/TurbulenceModel.pic}
\caption{\fort{TurbulenceModel\_t} Data Structure}
\label{f:TurbulenceModel}
\end{figure}

\clearpage
\begin{figure}[!htp]
\centering
\input{figures/ThermalRelaxationModel.pic}
\caption{\fort{ThermalRelaxationModel\_t} Data Structure}
\label{f:ThermalRelaxationModel}
\end{figure}

\clearpage
\begin{figure}[!htp]
\centering
\input{figures/ChemicalKineticsModel.pic}
\caption{\fort{ChemicalKineticsModel\_t} Data Structure}
\label{f:ChemicalKineticsModel}
\end{figure}

\clearpage
\begin{figure}[!htp]
\centering
\input{figures/EMElectricFieldModel.pic}
\caption{\fort{EMElectricFieldModel\_t} Data Structure}
\label{f:EMElectricFieldModel}
\end{figure}

\clearpage
\begin{figure}[!htp]
\centering
\input{figures/EMMagneticFieldModel.pic}
\caption{\fort{EMMagneticFieldModel\_t} Data Structure}
\label{f:EMMagneticFieldModel}
\end{figure}

\clearpage
\begin{figure}[!htp]
\centering
\input{figures/EMConductivityModel.pic}
\caption{\fort{EMConductivityModel\_t} Data Structure}
\label{f:EMConductivityModel}
\end{figure}

\clearpage
\begin{figure}[!htp]
\centering
\input{figures/ConvergenceHistory.pic}
\caption{\fort{ConvergenceHistory\_t} Data Structure}
\label{f:ConvergenceHistory}
\end{figure}

\clearpage
\begin{figure}[!htp]
\centering
\input{figures/IntegralData.pic}
\caption{\fort{IntegralData\_t} Data Structure}
\label{f:IntegralData}
\end{figure}

\clearpage
\begin{figure}[!htp]
\centering
\input{figures/ReferenceState.pic}
\caption{\fort{ReferenceState\_t} Data Structure}
\label{f:ReferenceState}
\end{figure}

\clearpage
\begin{figure}[!htp]
\centering
\input{figures/DataArray.pic}
\caption{\fort{DataArray\_t} Data Structure}
\label{f:DataArray}
\end{figure}

\clearpage
\begin{figure}[!htp]
\centering
\input{figures/DimensionalUnits.pic}
\caption{\fort{DimensionalUnits\_t} Data Structure}
\label{f:DimensionalUnits}
\end{figure}

\clearpage
\begin{figure}[!htp]
\centering
\input{figures/DimensionalExponents.pic}
\caption{\fort{DimensionalExponents\_t} Data Structure}
\label{f:DimensionalExponents}
\end{figure}

\clearpage
\begin{figure}[!htp]
\centering
\input{figures/Family.pic}
\caption{\fort{Family\_t} Data Structure}
\label{f:Family}
\end{figure}

\clearpage
\begin{figure}[!htp]
\centering
\input{figures/FamilyBC.pic}
\caption{\fort{FamilyBC\_t} Data Structure}
\label{f:FamilyBC}
\end{figure}

\clearpage
\begin{figure}[!htp]
\centering
\input{figures/GeometryReference.pic}
\caption{\fort{GeometryReference\_t} Data Structure}
\label{f:GeometryReference}
\end{figure}

\clearpage
\begin{figure}[!htp]
\centering
\input{figures/BaseIterativeData.pic}
\caption{\fort{BaseIterativeData\_t} Data Structure}
\label{f:BaseIterativeData}
\end{figure}

\clearpage
\begin{figure}[!htp]
\centering
\input{figures/ZoneIterativeData.pic}
\caption{\fort{ZoneIterativeData\_t} Data Structure}
\label{f:ZoneIterativeData}
\end{figure}

\clearpage
\begin{figure}[!htp]
\centering
\input{figures/RigidGridMotion.pic}
\caption{\fort{RigidGridMotion\_t} Data Structure}
\label{f:RigidGridMotion}
\end{figure}

\clearpage
\begin{figure}[!htp]
\centering
\input{figures/ArbitraryGridMotion.pic}
\caption{\fort{ArbitraryGridMotion\_t} Data Structure}
\label{f:ArbitraryGridMotion}
\end{figure}

\clearpage
\begin{figure}[!htp]
\centering
\resizebox{!}{\textheight-2\baselineskip}{\input{figures/UserDefinedData.pic}}
%\input{figures/UserDefinedData.pic}
\caption{\fort{UserDefinedData\_t} Data Structure}
\label{f:UserDefinedData}
\end{figure}

\clearpage
\begin{figure}[!htp]
\centering
\input{figures/Gravity.pic}
\caption{\fort{Gravity\_t} Data Structure}
\label{f:Gravity}
\end{figure}


\end{document}
